\documentclass[11pt]{article}

\usepackage[margin=1in]{geometry}
\usepackage{amsmath,amssymb,amsthm,mathtools}
\usepackage{booktabs}
\usepackage[hidelinks]{hyperref}
\usepackage{xcolor}

\newtheorem{theorem}{Theorem}
\newtheorem{definition}{Definition}
\newtheorem{claim}{Claim}
\newtheorem{assumption}{Assumption}
\newtheorem{corollary}{Corollary}
\theoremstyle{remark}
\newtheorem{remark}{Remark}
\setlength{\emergencystretch}{3em}

\newcommand{\R}{\mathbb{R}}
\newcommand{\PMF}{\operatorname{PMF}}
\newcommand{\Ent}{\operatorname{H}}
\newcommand{\Novelty}{\operatorname{N}}
\newcommand{\RewardB}{R_B}
\newcommand{\RewardC}{R_C}
\newcommand{\RewardD}{R_D}
\newcommand{\RewardDT}{\tilde{R}_D}
\newcommand{\RewardN}{R_N}
\newcommand{\RewardTrain}{R_{\mathrm{train}}}
\newcommand{\RewardFull}{R_{\mathrm{full}}}
\newcommand{\Feas}{\operatorname{Feas}}
\newcommand{\E}{\mathbb{E}}
\newcommand{\LeanName}[1]{\texttt{\detokenize{#1}}}

% Callout box for protocol-critical warnings (Goodhart enforcement).
\newcommand{\calloutbox}[2]{%
  \begin{center}
  \fcolorbox{black}{gray!10}{\parbox{0.92\linewidth}{%
    \textbf{#1.}\ #2}}
  \end{center}}

\title{Formal Creativity for MAS-LLMs:\\A Lean-Checked Bridge from Information Objectives to Benchmark Claims}
\author{}
\date{}

\begin{document}
\maketitle

\begin{abstract}
We introduce Creative Update Efficiency (CUE), a receiver-grounded creativity
metric defined as Brier-score delta per output bit: an output is credited only
for the decision-quality improvement it induces in a receiver agent,
normalized by its length in bits. CUE is proved to be a lower bound on
normalized Shannon Value of Information
(\texttt{CUE\_voi\_equivalence}) and bounded by the Shannon maximum
$\log|\mathcal{Z}|/|y|_{\mathrm{bits}}$
(\texttt{shannon\_upper\_bound\_normalization}), where the maximum-entropy
bound is itself derived from the Gibbs inequality rather than assumed.
The reported \textbf{Benchmark Score} is defined by
\[
  R_{\mathrm{creativity}}(q,y,s)
  =
  1\bigl[\mathrm{CUE}(y,A,T)>0\bigr]
  \cdot 1\bigl[R_D(y,\mathcal{P})>\delta_D\bigr]
  \cdot\Bigl(
    \mathrm{CUE}(y,A,T)\cdot\bigl(1+\alpha\,R_B^{\to A}(y,A)\bigr)
    +\lambda_G\,G_k(q,y)
  \Bigr)
\]
(\LeanName{Rcreativity}): both $1[\mathrm{CUE}>0]$ and
$1[R_D>\delta_D]$ are multiplicative safety floors on the
\emph{entire} score --- including $\lambda_G G_k$ --- so neither CUE failure
nor D-gate failure leaves an additive escape hatch. There is no
compression-novelty term $R_C$. Boundedness, monotonicity in CUE, and
total-score gate closure are machine-checked. Trajectory diagnostics (Step-CUE, Diverge-Converge)
and a multi-agent bridge ($G_k\Rightarrow$ joint CUE improvement) accompany
the score. A Dual-Process Separation Principle is proved as a structural
theorem: the training signal $R_{\mathrm{train}}=R_B+\tilde\lambda_D\tilde R_D$
depends only on live resampled quantities, so its gradient is independent of
every frozen benchmark object, and a Structural Goodhart Collapse theorem shows
that violating any protocol condition admits a rank reversal that refutes the
proxy-faithfulness guarantee. The Lean~4 development certifies these claims;
remaining assumptions are centralized in a Known Gaps Registry. The document
separates frozen benchmark objectives (Part~\ref{part:benchmark}) from the
2-CGRPO training signal (Part~\ref{part:training}); machine-checked
Goodhart-resistance conditions and a monitorable proxy faithfulness invariant
(Part~\ref{part:proxy}) govern when training results may be compared against
the benchmark.
\end{abstract}

\section{The Problem with Scalar Creativity Evaluation}

Current MAS-LLM evaluation often treats creativity as a scalar quality score:
a human or LLM judge rates an output against a rubric. This is insufficient for
three reasons.

First, it is reductive. A scalar score collapses a multi-dimensional property
into a single number and hides the distinction between novelty, usefulness,
diversity, surprise, downstream impact, and interaction structure.

Second, it is evaluator-dependent. LLM judges can over-rate LLM-generated ideas
relative to expert human evaluation, producing a novelty mirage: outputs appear
creative to the evaluator because they are fluent or stylistically unusual, not
because they introduce new information.

Third, scalar evaluation misses the point of multi-agent systems. Creativity in
a collective system is not merely the sum of individual creative outputs. The
central object is what emerges from interaction that no individual agent could
produce alone.

\begin{claim}[Core thesis]
Creativity in MAS-LLMs should be defined as a property of the system's
information dynamics, not as a property of individual outputs.
\end{claim}

\subsection*{Framework Preview}

We define creativity as utility-gated informational contribution. The
feasibility condition is
\[
  \Feas(q,y) \iff U(q,y)\ge \tau.
\]
Within the feasible set, B measures downstream expansiveness and D measures
probe-relative structural influence on a fixed hold-out probe set
$\mathcal{P}$ (defined by the protocol in
Appendix~\ref{app:probe-protocol}). The sole reported score is the
\textbf{Benchmark Score}
\begin{equation}
\label{eq:Rcreativity}
  R_{\mathrm{creativity}}(q,y,s)
  :=
  1\bigl[\mathrm{CUE}(y,A,T)>0\bigr]
  \cdot 1\bigl[\RewardD(y,\mathcal{P})>\delta_D\bigr]
  \cdot\Bigl(
    \mathrm{CUE}(y,A,T)\cdot\bigl(1+\alpha\,R_B^{\to A}(y,A)\bigr)
    +\lambda_G\,G_k(q,y)
  \Bigr),
\end{equation}
formalized as \LeanName{Rcreativity}
(Section~\ref{sec:benchmark-score}): \LeanName{cueGate} and
\LeanName{dGate} multiply the whole bracket
$\mathrm{CUE}\cdot(1+\alpha R_B^{\to A})+\lambda_G G_k$, enforcing a
strict safety floor (either gate failure zeroes the score, including
interaction gain). There is no compression-novelty component $R_C$. An
optional additive B--D diagnostic $\RewardB+\lambda_D\RewardD$
(\LeanName{benchmarkReward}) supports rank-preservation correspondence and
is \emph{not} $R_{\mathrm{creativity}}$.

\calloutbox{This is not the training objective}{$\RewardD$ requires a frozen probe set $\mathcal{P}$ and must never receive gradient updates. The training objective is $\RewardTrain=R_B+\tilde\lambda_D\tilde R_D$, defined in Part~\ref{part:training}.}

\subsection*{Contributions Hierarchy}

\calloutbox{Contributions, in order of priority}{%
\textbf{1. Primary — CUE as VOI-grounded multiplicative base (new).} Creative
Update Efficiency, receiver decision-quality gain per output bit, proved to
lower-bound normalized Shannon Value of Information
(\LeanName{CUE_voi_equivalence}) and to respect the Shannon ceiling
(\LeanName{shannon_upper_bound_normalization}); the reported Benchmark Score
is~\eqref{eq:Rcreativity} (\LeanName{Rcreativity}).
\textbf{2. Secondary — Dual-Process Separation Principle.} Goodhart
resistance as an \emph{architectural theorem}
(\LeanName{dual_process_separation_principle}): the training gradient is
structurally independent of frozen benchmark objects, paired with the
Structural Goodhart Collapse theorem
(\LeanName{structural_goodhart_collapse}) showing each protocol violation
admits a refuting rank reversal.
\textbf{3. Tertiary — Trajectory diagnostics.} The Step-CUE curve is
monotone under the per-step gate
(\LeanName{CUE_trajectory_monotone_nondecreasing}), its saturation rate
$\gamma$ is positive exactly under continuous discovery
(\LeanName{step_CUE_gamma_positivity_characterization}), and the
Diverge-Converge score is achievable and non-trivial
(\LeanName{DC_score_existence}, \LeanName{monotone_increasing_not_DC}).
\textbf{4. Supporting — B/D information-theoretic scaffolding.} Entropy
monotonicity, mixture anti-collapse (now axiom-free), utility gating,
probe-relative deformation, baseline contrast, estimator rank preservation,
and the multi-agent interaction-gain bridge.}

\subsection*{Relationship to ComplLLM}

ComplLLM (Guo et al., 2026) post-trains a decision-assistant LLM using
complementary information as reward: an output scores only if it complements
existing agent decisions, validated in verifiable decision-support domains.
CUE inherits from ComplLLM the receiver-outcome focus and the use of decision
quality as a training signal. It generalizes ComplLLM in three ways: (a) it
applies to open-ended creative tasks without binary payoff decompositions,
because the Brier-delta is defined against any finite decision-relevant state
space; (b) it normalizes by output length in bits and connects the normalized
quantity to a Shannon upper bound (\LeanName{shannon_upper_bound_normalization});
and (c) it is embedded in a dual-process Goodhart-resistance protocol that
separates frozen benchmark objects from live training proxies
(Part~\ref{part:dps}).

\subsection*{English Roadmap of the Lean Argument}

The Lean development supports this thesis by proving that the proposed
information-dynamic objectives behave exactly as the paper claims. The proof
strategy is simple in spirit. First, define a reward B that gives credit to
outputs only when they pass a utility gate, and then ranks feasible outputs by
the entropy of the downstream states they induce. Lean proves that once the
gate is open, higher downstream entropy strictly means higher B reward. It also
proves that a lower-entropy feasible output cannot be optimal when a
higher-entropy feasible output exists. If the high-entropy output receives
sufficient probability mass, mixing over feasible outputs cannot do worse than
collapsing to a strictly dominated lower-entropy output. In ordinary terms: B
formally rules out the kind of collapse where a system keeps choosing a safe,
narrow, predictable answer despite the availability of a more generative
feasible alternative.

Second, define a reward C that gives credit to outputs only when they pass the
same utility gate, and then ranks them by how costly they are to encode relative
to an existing corpus. Lean proves that C gives zero reward to infeasible
outputs no matter how novel they look, that exact copies are bounded by the
compressor's copy-minimality axiom, and that padding with zero-information
content cannot improve the novelty score. In ordinary terms: C separates
genuine novelty from three common failure modes: nonsense, plagiarism, and
length padding. This is exploratory creativity in Boden's sense: novelty
\emph{within} a fixed conceptual space.

Third, for empirical implementation, define surface novelty N as KL divergence
between an output's n-gram PMF and the corpus PMF. N is not a fourth creativity
dimension; it is the tractable surface approximation to C. Lean proves its gate,
nonnegativity, and cross-entropy link in the computational approximations
appendix.

Fourth, define D as probe-relative structural influence: corpus deformation gain
on a fixed finite probe set $\mathcal{P}$. Under \LeanName{ProbeCompressor}, Lean
proves that D is utility-gated, order-preserving, nonnegative, and zero for exact
corpus copies, for any fixed $\mathcal{P}$. The benchmark operationalizes
$\mathcal{P}$ by the hold-out protocol in Appendix~\ref{app:probe-protocol}.
High $D$ correlates with transformational creativity when $\mathcal{P}$ is
domain-representative; the formal theorems do not identify $D$ with Boden--Wiggins
transformation outright. High C measures departure from the current corpus; high
$D$ measures how much $s$ changes compression of the frozen probes.

Fifth, Lean proves that these objectives add something real beyond obvious
baselines. A utility-only reward cannot distinguish two feasible outputs with
the same utility but different downstream entropy; B can. A novelty-only reward
can assign positive score to an infeasible string; C cannot. D adds a separate
probe-relative invariant: it measures whether incorporating an output changes
compression of a frozen hold-out probe set, while Lean proves the score is
utility-gated, order-preserving, nonnegative, and zero for exact corpus copies.
These contrast theorems are important because they show that B, C, and D are not
just renamings of quality or novelty rubrics; the surface-N proxy inherits the
same gate and nonnegativity discipline in the appendix. They enforce different
invariants.

Sixth, Lean bridges the abstract objectives to measurement and training. If a
benchmark estimator has uniformly bounded error, then entropy and novelty
rankings are preserved whenever the true gap is larger than twice the error
bound. This gives the benchmark an explicit resolution condition: when measured
differences exceed $2\epsilon$, the ranking is not an artifact of estimator
noise. Lean also proves that a simplified DPO update constructed from a
B-valid preference pair increases the winner's log-odds over the loser, so the
training signal moves in the direction certified by the objective. The additive
\LeanName{benchmarkReward} diagnostic combines B and D for group-relative RL
correspondence; the reported score remains~\eqref{eq:Rcreativity}.

Finally, Lean extends the argument to multi-agent settings. It defines an
interaction gain: the difference between joint downstream entropy and the best
individual downstream entropy. If this gain is positive, Lean proves that
neither individual agent alone achieves the joint entropy. The gain is zero
exactly in the degenerate case where the joint entropy equals the best
individual entropy (\LeanName{gain_zero_iff_joint_eq_max}); notably,
independent product kernels with positive marginals have strictly positive
gain by entropy additivity, which is why the framework pairs the gain with the
faithfulness assumption and the receiver-level CUE bridge
(\LeanName{Gk_implies_CUE_improvement}) rather than treating positive gain
alone as synergy. If the gain
exceeds twice the benchmark error, Lean proves that the estimator detects this
irreducibility.
The $k$-agent extension and O-information sign bridge are the general versions;
the two-agent case is the readable base case.

Taken together, the Lean proofs establish the paper's central logical chain:
utility-gated entropy prevents collapse; utility-gated compression novelty
prevents nonsense, copying, and padding; bounded estimators preserve the
theoretical rankings; DPO-style training moves in the certified direction; and
faithful non-product multi-agent kernels provide a formal proxy for emergent
coordination.


%======================================================================
\part{Receiver-Grounded Creativity Foundation}
\label{part:cue}
%======================================================================

\noindent
This part is the definitional anchor of the framework. Creativity is grounded
in the \emph{receiver}: an output is creative only insofar as it improves a
receiver agent's downstream decisions, efficiently per bit. Everything else in
the document --- the B/C/D components, the trajectory diagnostics, the
multi-agent gain, the training signal --- is subordinate to this gate.

\section{Creative Update Efficiency (CUE)}

\begin{definition}[CUE model]
\LeanName{CUEModel}. A receiver-grounded CUE model over a finite
decision-relevant state space $\mathcal{Z}$ carries the receiver's posterior
belief after the output, the Brier-score improvement
$\Delta_{\mathrm{Brier}}(y) \ge 0$ from prior to posterior, the Shannon value
of information $\mathrm{VOI}(y)$ of the output for the receiver's decision,
and the output length $|y|_{\mathrm{bits}} > 0$. Two structural facts are
carried as certified fields: the proper-scoring-rule bound
$\Delta_{\mathrm{Brier}} \le \mathrm{VOI}$, and the residual-entropy budget
$\mathrm{VOI} \le \Ent(\mathrm{posterior})$.
\end{definition}

\begin{definition}[Creative Update Efficiency]
\LeanName{cue}. For a CUE model $m$,
\[
  \mathrm{CUE}(y, A, T)
  :=
  \frac{\Delta_{\mathrm{Brier}}(y)}{|y|_{\mathrm{bits}}},
\]
the receiver decision-quality gain per output bit. Nonnegativity is
\LeanName{cue_nonneg}.
\end{definition}

\begin{theorem}[PROOF-01: CUE lower-bounds normalized VOI]
\LeanName{CUE_voi_equivalence}. For any CUE model,
\[
  \mathrm{CUE}(y, A, T) \le \frac{\mathrm{VOI}(y, A, T)}{|y|_{\mathrm{bits}}},
\]
with equality when the Brier improvement saturates its proper-scoring-rule
bound --- the calibrated binary case (\LeanName{CUE_voi_equality}).
\end{theorem}
\begin{proof}
Divide the proper-scoring-rule bound $\Delta_{\mathrm{Brier}} \le \mathrm{VOI}$
by the positive bit length; Lean closes with
\texttt{div\_le\_div\_iff\_of\_pos\_right}. Equality substitutes
$\Delta_{\mathrm{Brier}} = \mathrm{VOI}$.
\end{proof}

\begin{theorem}[Maximum entropy on a finite alphabet]
\LeanName{maxEntropy_finite_alphabet}. For any PMF $p$ over a finite nonempty
alphabet $\mathcal{Z}$,
\[
  \Ent(p) \le \log|\mathcal{Z}|.
\]
\end{theorem}
\begin{proof}
No new axiom is introduced: the KL divergence of $p$ against the uniform
distribution evaluates to $\log|\mathcal{Z}| - \Ent(p)$
(\LeanName{klDivergence_uniform_eq}, using the finite mass identity
\LeanName{sum_toReal_eq_one}), and the existing Gibbs inequality
\LeanName{klDivergence_nonneg} makes it nonnegative.
\end{proof}

\begin{theorem}[PROOF-02: Shannon upper bound normalization]
\LeanName{shannon_upper_bound_normalization}. For any CUE model over
$\mathcal{Z}$,
\[
  \mathrm{CUE}(y, A, T) \le \frac{\log|\mathcal{Z}|}{|y|_{\mathrm{bits}}}.
\]
\end{theorem}
\begin{proof}
Chain $\Delta_{\mathrm{Brier}} \le \mathrm{VOI} \le \Ent(\mathrm{posterior})
\le \log|\mathcal{Z}|$, the last step by
\LeanName{maxEntropy_finite_alphabet}, and divide by the positive bit length.
\end{proof}

\begin{theorem}[CUE primary gate]
\LeanName{CUE_primary_gate}. Zero or negative CUE disqualifies all downstream
metrics: if $\mathrm{CUE}(y,A,T) \not> 0$, the gated creativity score of $y$
is zero regardless of its B, C, or D values (\LeanName{cueGatedScore}).
\end{theorem}

\calloutbox{CUE gates everything}{The B/C/D components of
Part~\ref{part:definition} are \emph{components of the CUE-gated score}, not
co-equal creativity dimensions. An output that fails the CUE gate receives no
creativity credit no matter how expansive, novel, or structurally influential
it is.}

\section{Receiver Expansion vs.\ Clarification}

The bidirectional receiver label $R_B^{\to A}$ records whether an output
\emph{expands} the receiver's downstream generative space (high receiver
entropy) or \emph{clarifies} it (low receiver entropy).

\begin{theorem}[PROOF-06: expansion/clarification separation]
\LeanName{receiver_expansion_clarification_separation}. The label
$R_B^{\to A}$ is not a monotone function of CUE: there exist outputs
$y_1, y_2$ with equal CUE but opposite labels --- one expanding
($R_B^{\to A}(y_1) > \delta_E$) and one clarifying
($R_B^{\to A}(y_2) < \delta_E$).
\end{theorem}
\begin{proof}
Explicit finite construction mirroring \LeanName{BC_tension_example}: two
labeled outputs (\LeanName{LabeledOutput}) share the same underlying CUE model
on $\operatorname{Fin}(2)$ (hence genuinely equal CUE), while their receiver
expansion magnitudes are $1$ (maximal, uniform-generation scale) and $0$
(point-mass commitment), with threshold $\delta_E = 1/2$ strictly between.
\end{proof}

\section{The Benchmark Score: CUE as Multiplier}
\label{sec:benchmark-score}
\label{sec:shifted-cue}

\paragraph{Canonical definition.}
The reported Benchmark Score is exactly~\eqref{eq:Rcreativity}:
\[
  R_{\mathrm{creativity}}(q,y,s)
  =
  \mathbb{1}\bigl[\mathrm{CUE}(y,A,T)>0\bigr]
  \cdot\mathbb{1}\bigl[R_D(y,\mathcal{P})>\delta_D\bigr]
  \cdot\Bigl(
    \mathrm{CUE}(y,A,T)\cdot\bigl(1+\alpha\,R_B^{\to A}(y,A)\bigr)
    +\lambda_G\,G_k(q,y)
  \Bigr).
\]
Lean name: \LeanName{Rcreativity} (abbrev.\ \LeanName{benchmarkScore}).
Gates are \LeanName{cueGate} and \LeanName{dGate}. There is no $R_C$ term.
Placing $\lambda_G G_k$ \emph{inside} both indicators closes the additive
escape hatch that an ungated interaction-gain term would create: any CUE or
$R_D$ violation immediately nullifies the total reward (multiplicative safety
floor). Proof obligations of this regime are discharged as follows.

\begin{theorem}[CUE-gate closure --- total zero]
\LeanName{Rcreativity_cue_gate_closed}. If $\mathrm{CUE}\le 0$, then
$R_{\mathrm{creativity}}=0$ (including the $\lambda_G G_k$ term).
\end{theorem}

\begin{theorem}[D-gate closure --- total zero]
\LeanName{Rcreativity_d_gate_closed}. If $R_D\le\delta_D$, then
$R_{\mathrm{creativity}}=0$ (including the $\lambda_G G_k$ term). The
magnitude of $R_D$ never enters the score (\LeanName{D2_dominance} is
diagnostic-only here).
\end{theorem}

\begin{theorem}[Nonnegativity]
\LeanName{Rcreativity_nonneg}. Under $\alpha>0$, $r_b\in[0,1]$,
$0\le\lambda_G$, and $0\le G_k$, one has $0\le R_{\mathrm{creativity}}$.
\end{theorem}

\begin{theorem}[CUE-as-multiplier ceiling]
\LeanName{Rcreativity_bounded}.
$R_{\mathrm{creativity}}\le
(\log|Z|/|y|_{\mathrm{bits}})\cdot(1+\alpha)+\lambda_G G_k$,
via \LeanName{shannon_upper_bound_normalization} and
\LeanName{receiver_expansion_bounded}.
\end{theorem}

\begin{theorem}[Monotonicity in CUE]
\LeanName{Rcreativity_monotone_in_CUE}. With both gates open and fixed
$(\alpha,r_b,\lambda_G,G_k)$, strictly larger CUE yields a strictly larger
score: CUE's magnitude is load-bearing.
\end{theorem}

\begin{theorem}[Positivity matches CUE]
\LeanName{Rcreativity_pos_iff_CUE_pos}. With the D-gate open and
$\lambda_G G_k=0$, one has $0<R_{\mathrm{creativity}}$ iff $0<\mathrm{CUE}$.
\end{theorem}

\begin{theorem}[Unit reconciliation]
\LeanName{unit_conversion_sign_and_order_invariant}. Positive rescaling
preserves sign and order. CUE is bits-normalized while $R_B^{\to A}$ is
nats-based; the residual scale is absorbed by the free-fit hyperparameter
$\alpha$ (KG-6), explicitly weaker than a proved $\lambda$-normalization.
\end{theorem}

\paragraph{Proof flow (central claims).}
\begin{enumerate}
  \item \emph{CUE is a valid base:} PROOF-01/02 give VOI lower-bound and
  Shannon ceiling; \LeanName{cueGate} is a safety floor on the whole score.
  \item \emph{$R_D$ is a filter, not a summand:} \LeanName{dGate} likewise
  zeroes the entire score; magnitude is discarded
  (\LeanName{D2_dominance} diagnostic-only). The former $R_C$-coupled D3
  bound is replaced by \LeanName{D3_rewardD_bound_standalone}.
  \item \emph{Expansion bonus is controlled:}
  \LeanName{receiver_expansion_bounded} gives $r_b\in[0,1]$, so the factor
  $1+\alpha r_b$ lies in $[1,1+\alpha]$. Theorem~5's separate-label result
  (\LeanName{receiver_expansion_clarification_separation}) remains valid as
  a diagnostic statement; folding $R_B^{\to A}$ into~\eqref{eq:Rcreativity}
  is a deliberate, documented exception for the reported score.
  \item \emph{$G_k$ has no ungated path:} $\lambda_G G_k$ sits inside both
  indicators, so Theorems~6--7 prove total-score zero on either gate failure
  (strengthening the earlier ``reduces to $\lambda_G G_k$'' statements).
  \LeanName{Gk_implies_CUE_improvement} / \LeanName{Gk_gated_by_CUE_shifted}
  remain compatible with this floor.
  \item \emph{Former $R_C$ dependencies are bridged:}
  \LeanName{lambda_D_normalization_standalone} and
  \LeanName{paretoFrontier_nonempty_BD} replace $\lambda_C$- and
  three-way-C-coupled commensurability/Pareto claims for any residual B--D
  diagnostic.
\end{enumerate}

\paragraph{Optional shifted calibration.}
The scalar $\mathrm{CUE}_f:=\mathrm{CUE}-f$ (\LeanName{CUE_shifted}, KG-5)
and \LeanName{Rcreativity_shifted} are retained only as calibration variants;
they are \emph{not} the canonical definition~\eqref{eq:Rcreativity}.

\begin{theorem}[Shifted Shannon ceiling]
\LeanName{CUE_shifted_bound}. For every CUE model $m$ and threshold $f$,
$\mathrm{CUE}_f \le \log|Z| / |y|_{\mathrm{bits}} - f$.
\end{theorem}

\begin{theorem}[Universal expansion bound]
\LeanName{receiver_expansion_bounded}. With
$R_B^{\to A} = H(\mathrm{Gen}(A \mid y, T)) / \log|W|$
(\LeanName{receiverExpansionNorm}), one has $0 \le R_B^{\to A} \le 1$.
\end{theorem}

\begin{theorem}[$G_k$ gated by shifted CUE]
\LeanName{Gk_gated_by_CUE_shifted}. Under
\LeanName{Gk_implies_CUE_improvement}, if $G_k > 0$ and
$f < \mathrm{CUE}$, then $G_k \cdot \mathbb{1}[\mathrm{CUE}_f > 0] > 0$.
\end{theorem}

\section{Formal Correspondence Table}

\begin{center}
\scriptsize
\begin{tabular}{@{}p{0.22\linewidth}p{0.48\linewidth}p{0.24\linewidth}@{}}
\toprule
Proposal construct & Lean theorem & Status \\
\midrule
$\mathrm{CUE}(y,A,T)$ & \LeanName{cue}, \LeanName{CUE_voi_equivalence} & Machine-checked (PROOF-01) \\
Shannon normalization & \LeanName{shannon_upper_bound_normalization} & Machine-checked (PROOF-02) \\
\textbf{Benchmark Score} $R_{\mathrm{creativity}}$ & \LeanName{Rcreativity} & Canonical def.~\eqref{eq:Rcreativity} \\
CUE / D gates (total zero) & \LeanName{cueGate}, \LeanName{dGate}, \LeanName{Rcreativity_cue_gate_closed}, \LeanName{Rcreativity_d_gate_closed} & Score $=0$ if either fails \\
Boundedness / monotone & \LeanName{Rcreativity_bounded}, \LeanName{Rcreativity_monotone_in_CUE} & Machine-checked \\
Positivity & \LeanName{Rcreativity_pos_iff_CUE_pos} & Machine-checked \\
Unit reconciliation & \LeanName{unit_conversion_sign_and_order_invariant} & Machine-checked \\
$R_B^{\to A}$ bound & \LeanName{receiver_expansion_bounded} & Machine-checked \\
$G_k$ decision bridge & \LeanName{Gk_implies_CUE_improvement}, \LeanName{Gk_gated_by_CUE_shifted} & Machine-checked (PROOF-07) \\
D3 standalone & \LeanName{D3_rewardD_bound_standalone} & Bridge ($R_C$-free) \\
$\lambda_D$ standalone & \LeanName{lambda_D_normalization_standalone} & Bridge (B--D diagnostic) \\
B--D Pareto & \LeanName{paretoFrontier_nonempty_BD} & Bridge (B--D diagnostic) \\
DPS / Goodhart & \LeanName{dual_process_separation_principle}, \LeanName{structural_goodhart_collapse} & Machine-checked \\
\bottomrule
\end{tabular}
\end{center}

%======================================================================
\part{Creativity Definition (Measurement)}
\label{part:definition}
%======================================================================

\noindent
This part is organized in three tiers subordinate to the CUE foundation of
Part~\ref{part:cue}: \textbf{Tier A (multiplicative base)} is gated CUE;
\textbf{Tier B (structural filters)} comprises the $R_D$ probe-relative gate
and the $R_B^{\to A}$ expansion bonus; \textbf{Tier C (interaction
certification)} is the multi-agent gain $G_k$. The sole reported measurement is
the Benchmark Score~\eqref{eq:Rcreativity}
(\LeanName{Rcreativity}, Section~\ref{sec:benchmark-score}).
It is \emph{not} the quantity the trainer optimizes: the training signal
$\RewardTrain=\RewardB+\tilde\lambda_D\RewardDT$
(Part~\ref{part:training}) replaces frozen $\RewardD$ by the live proxy.
The additive intermediate $\RewardB+\lambda_D\RewardD$
(\LeanName{benchmarkReward}) serves B/D rank-preservation correspondence only
and is never identified with $R_{\mathrm{creativity}}$. Compression novelty
$R_C$ is not a score component; \LeanName{AbstractCompressor} remains solely
as infrastructure for \LeanName{ProbeCompressor}/$R_D$.

\section{Formal Framework}

\subsection{The B Reward}

The Lean structure \LeanName{BModel} contains a kernel
\[
  \texttt{kernel} : Q \to Y \to \PMF(Z),
\]
a utility function $U : Q \to Y \to \R$, and a threshold $\tau \in \R$.
The project defines
\[
  \Feas_m(q,y) \iff \tau \le U(q,y)
\]
and
\[
  \RewardB(q,y) :=
    \begin{cases}
      H_q(y), & \Feas_m(q,y),\\
      0, & \text{otherwise}.
    \end{cases}
\]

\subsection{Compressor infrastructure (not a score component)}

$R_C$ is \emph{not} part of $R_{\mathrm{creativity}}$. The Lean class
\LeanName{AbstractCompressor} is retained solely so that
\LeanName{ProbeCompressor} (and hence $R_D$) can use code-length /
novelty primitives:

\begin{enumerate}
  \item a code length function $L : \text{List}(\alpha) \to \R$;
  \item nonnegativity: $0 \le L(s)$ for all strings $s$;
  \item copy minimality:
  \[
    L(\mathcal{H} \mathbin{++} s) - L(\mathcal{H}) \le L(s).
  \]
\end{enumerate}
The novelty score is
\[
  N(s \mid \mathcal{H}) = L(\mathcal{H} \mathbin{++} s) - L(\mathcal{H}),
\]
and the gated reward is
\[
  \RewardC(s,\mathcal{H},u,\tau) :=
    \begin{cases}
      N(s \mid \mathcal{H}), & \tau \le u,\\
      0, & \text{otherwise}.
    \end{cases}
\]
This $\RewardC$ object is \emph{not} used in~\eqref{eq:Rcreativity}; it is
shown only because $N(\,\cdot\mid\mathcal{H})$ underpins $R_D$.
For padding results, \LeanName{NormalizedCompressor} adds:
\[
  L(\mathcal{H} \mathbin{++} (s \mathbin{++} p))
    \le L(\mathcal{H} \mathbin{++} s) + L(p).
\]

\subsection{The D Reward}

The Lean class \LeanName{ProbeCompressor} extends the compressor setting with a
fixed finite probe set $\mathcal{P}$. The deformation score measures how much
incorporating $s$ into the corpus reduces compression cost on those probes:
\[
  D(s,\mathcal{H},\mathcal{P})
  :=
  \sum_{s'\in\mathcal{P}}
  \bigl(N(s'\mid\mathcal{H})-N(s'\mid\mathcal{H}\cup\{s\})\bigr).
\]
The gated reward is
\[
  \RewardD(s,\mathcal{H},\mathcal{P},u,\tau) :=
    \begin{cases}
      D(s,\mathcal{H},\mathcal{P}), & \tau \le u,\\
      0, & \text{otherwise}.
    \end{cases}
\]
D therefore measures probe-relative structural influence: it is not novelty
relative to the current corpus, but the extent to which adding the output changes
compression of a frozen hold-out probe set.

\subsection{Combined Training Objective}

In a DPO-style training regime, preference pairs are built inside the feasible
set. A B-pair prefers the feasible output with higher downstream entropy; a
D-pair prefers the feasible output with higher probe-relative deformation gain.
The conceptual combined loss is
\[
  \mathcal{L}(\theta)
    = \mathcal{L}_{B\text{-DPO}}(\theta)
      + \lambda_D \mathcal{L}_{D\text{-DPO}}(\theta),
\]
with no $R_C$ / C-DPO channel. B governs downstream impact and D governs
structural influence on a fixed probe set. Reported evaluation uses
$R_{\mathrm{creativity}}$~\eqref{eq:Rcreativity}, not this training loss.

\subsection{Commensurability: B--D diagnostic normalization}

The optional additive diagnostic $\RewardB+\lambda_D\RewardD$ is justified
because $H_q(y)$ and $D(s,\mathcal{H},\mathcal{P})$ are Shannon quantities
in nats; $\lambda_D$ is a dimensionless scale. The reported Benchmark Score
is multiplicative~\eqref{eq:Rcreativity} and instead uses the unit-reconciliation
lemma of Section~\ref{sec:benchmark-score}. For the diagnostic:

\begin{theorem}[B ceiling]
\LeanName{rewardB_le_maxEntropy}. For every prompt and output,
$\RewardB(q,y) \le \log|\mathcal{Z}|$: gated-off outputs score
$0 \le \log|\mathcal{Z}|$ and feasible outputs score their downstream entropy,
bounded by \LeanName{maxEntropy_finite_alphabet} (Part~\ref{part:cue}).
\end{theorem}

\begin{theorem}[$\lambda_D$ normalization]
\LeanName{lambda_D_normalization_standalone}. With
$\lambda_D=\log|\mathcal{Z}|/(|P|\cdot\log|\Sigma|\cdot\E[|s'|])$, where
$|P|$ is the probe set size and $\E[|s'|]$ is the mean probe string length,
\[
  \lambda_D \cdot \bigl(|P|\cdot\log|\Sigma|\cdot\E[|s'|]\bigr)
  = \log|\mathcal{Z}|,
\]
so the maximal $\lambda_D$-scaled D contribution coincides exactly with the
proved B ceiling.
\end{theorem}

$R_B$ and $R_D$ are two projections of the same generative event: the
agent's output simultaneously induces a downstream state distribution
(measured by B) and deforms the compression of domain-representative probes
(measured by D); normalization ensures neither projection dominates by
default in the diagnostic.

\subsection{Summary Criterion}

In this framework, a MAS-LLM behavior is creative when it scores positively
on~\eqref{eq:Rcreativity}: both CUE and $R_D$ gates open (else the score is
identically zero), with CUE magnitude, expansion bonus, and (when applicable)
interaction gain $G_k$ contributing inside that floor. High D on a
domain-representative $\mathcal{P}$ is an empirical correlate of
transformational creativity rather than transformation by definition.

\section{Lean-Checked Objective Guarantees}

This section records the exact mathematical assumptions under which the Lean
project proves each objective property. The proof checker verifies the
statements listed by their Lean names.

\subsection{Entropy Layer}

\begin{theorem}[Shannon entropy nonnegativity]
\LeanName{shannonEntropy_nonneg}. For any finite type $\alpha$ and any
$p : \PMF(\alpha)$,
\[
  0 \le -\sum_{x:\alpha} p(x) \log p(x).
\]
\end{theorem}
\begin{proof}
For each $x$, $p(x)$ is a probability mass, so $0\le p(x)\le 1$. Hence
$\log p(x)\le 0$ and therefore $p(x)\log p(x)\le 0$. Summing over the finite
type gives
\[
  \sum_x p(x)\log p(x)\le 0.
\]
Multiplying by $-1$ yields
\[
  0\le -\sum_x p(x)\log p(x).
\]
In Lean this is proved by applying \texttt{Real.mul\_log\_nonpos} pointwise
and then \texttt{Finset.sum\_nonpos}.
\end{proof}

\begin{theorem}[Point mass entropy]
\LeanName{shannonEntropy_point_mass}. For any finite type $\alpha$ and
$x_0:\alpha$,
\[
  H(\delta_{x_0}) = 0.
\]
\end{theorem}
\begin{proof}
For the point mass $\delta_{x_0}$, the probability is $1$ at $x_0$ and $0$ at
every other point. Thus every entropy summand is either $1\log 1$ or
$0\log 0$, both of which evaluate to $0$ in the real-valued definition used by
the project. Therefore the finite sum is zero. Lean proves this by splitting on
whether $x=x_0$ and simplifying \texttt{PMF.pure\_apply}.
\end{proof}

\begin{theorem}[Entropy order unfolding]
\LeanName{shannonEntropy_lt_iff}. For finite $\alpha$ and $p,q:\PMF(\alpha)$,
\[
  H(p)<H(q)
  \quad\Longleftrightarrow\quad
  -\sum_x p(x)\log p(x)<-\sum_x q(x)\log q(x).
\]
\end{theorem}
\begin{proof}
This is definitional: \LeanName{shannonEntropy} is exactly the finite negative
sum on both sides. Lean proves the equivalence by simplifying with the
definition of \LeanName{shannonEntropy}.
\end{proof}

\subsection{Lean Helper Lemmas for B}

The B model file also proves four helper lemmas used by later theorems.

\begin{theorem}[Reward on feasible outputs]
\LeanName{rewardB_of_feasible}. If $\Feas_m(q,y)$, then
\[
  \RewardB(q,y)=H_q(y).
\]
\end{theorem}
\begin{proof}
The definition of $\RewardB$ is an if-then-else on feasibility. Under the
assumption $\Feas_m(q,y)$, the true branch is selected. Lean proves this by
simplifying \texttt{rewardB}.
\end{proof}

\begin{theorem}[Reward on infeasible outputs]
\LeanName{rewardB_of_infeasible}. If $\neg\Feas_m(q,y)$, then
\[
  \RewardB(q,y)=0.
\]
\end{theorem}
\begin{proof}
Unfold the same if-then-else defining $\RewardB$. Under infeasibility, the false
branch is selected, which is $0$. Lean proves this by simplification.
\end{proof}

\begin{theorem}[B reward nonnegativity]
\LeanName{rewardB_nonneg}. For every $q,y$,
\[
  0\le \RewardB(q,y).
\]
\end{theorem}
\begin{proof}
If $y$ is feasible, then $\RewardB(q,y)=H_q(y)$, which is nonnegative by
\LeanName{shannonEntropy_nonneg}. If $y$ is infeasible, then
$\RewardB(q,y)=0$. Lean proves this by case-splitting on feasibility and
simplifying.
\end{proof}

\begin{theorem}[Feasible-set membership]
\LeanName{mem_feasibleSet_iff}. For finite $Y$,
\[
  y\in \operatorname{feasibleSet}(q)
  \quad\Longleftrightarrow\quad
  \Feas_m(q,y).
\]
\end{theorem}
\begin{proof}
The feasible set is defined as the finite universe filtered by the feasibility
predicate. Membership in this filtered universe is therefore equivalent to
feasibility itself. Lean proves this by simplifying \texttt{feasibleSet}.
\end{proof}

\subsection{B: Anti-Collapse and Entropic Ordering}
\label{sec:b-anti-collapse}

\begin{theorem}[B1: dominance]
\LeanName{B1_dominance}. Assume $Z$ is finite, $m : \texttt{BModel}(Q,Y,Z)$,
$q:Q$, and $y_1,y_2:Y$. If
\[
  \Feas_m(q,y_1), \qquad \Feas_m(q,y_2),
  \qquad H_q(y_2) < H_q(y_1),
\]
then
\[
  \RewardB(q,y_2) < \RewardB(q,y_1).
\]
\end{theorem}
\begin{proof}
Because both outputs are feasible, the utility gate is open for both:
\[
  \RewardB(q,y_i)=H_q(y_i)\qquad (i=1,2).
\]
Substituting these equalities into the desired conclusion reduces it exactly to
the hypothesis $H_q(y_2)<H_q(y_1)$. Lean proves this by simplifying the
definition of \texttt{rewardB} under the feasibility hypotheses.
\end{proof}

\noindent
Consequence: inside the feasible set, B rewards exactly the downstream entropy
ordering. The result is order-theoretic: once the gate is open, the reward
unfolds to entropy.

\begin{theorem}[B2: no entropy-dominated optimum]
\LeanName{B2_no_collapse}. Assume $Y$ and $Z$ are finite. Let
$m : \texttt{BModel}(Q,Y,Z)$, $q:Q$, and $y,y^\star:Y$. If
\[
  \Feas_m(q,y), \qquad \Feas_m(q,y^\star), \qquad
  H_q(y) < H_q(y^\star),
\]
and the feasible set $\{y' : \Feas_m(q,y')\}$ is nonempty, then
\[
  \RewardB(q,y)
    < \sup_{y' \in \Feas_m(q,\cdot)} \RewardB(q,y').
\]
\end{theorem}
\begin{proof}
By B1, the hypotheses imply
\[
  \RewardB(q,y)<\RewardB(q,y^\star).
\]
Since $y^\star$ is feasible, it belongs to the feasible finite set. A finite
supremum is strictly above a value whenever some member of the set is strictly
above that value. Applying this finite-supremum fact to $y^\star$ gives
\[
  \RewardB(q,y)
  <
  \sup_{y'\in\Feas_m(q,\cdot)}\RewardB(q,y').
\]
Lean uses \texttt{Finset.lt\_sup'\_iff} with witness $y^\star$.
\end{proof}

\noindent
Consequence: a deterministic policy concentrated on a strictly
entropy-dominated feasible output cannot be optimal for B.

\subsubsection*{Mixture rewards (B3)}

\noindent\textbf{Proof status (PROOF-09, resolved).}
The distributional Jensen extension \LeanName{B3b_jensen_anti_collapse}
formerly depended on \LeanName{entropy_mixture_concavity_axiom}. That axiom
has been \emph{removed and promoted to a theorem},
\LeanName{entropy_mixture_concavity}, proved from primitives via the concave
Jensen inequality (\texttt{ConcaveOn.le\_map\_sum}) applied to
\texttt{Real.concaveOn\_negMulLog}, with the pointwise mixture-kernel
evaluation \LeanName{mixtureKernel_apply_toReal} and the finite mass identity
\LeanName{sum_toReal_eq_one}. B3b is now a full theorem with no axiom
dependence.

\begin{theorem}[B3: mixture lower bound]
\LeanName{B3_mixture_lb}. Assume $Y$ and $Z$ are finite. Let $\pi:\PMF(Y)$.
If the support of $\pi$ is feasible, $y_{\mathrm{low}}$ and
$y_{\mathrm{high}}$ are feasible, $H_q(y_{\mathrm{low}})<H_q(y_{\mathrm{high}})$,
$\pi(y_{\mathrm{high}})>0$, and
\[
  \RewardB(q,y_{\mathrm{low}})
  \le
  \pi(y_{\mathrm{high}})\RewardB(q,y_{\mathrm{high}}),
\]
then
\[
  \RewardB(q,y_{\mathrm{low}})
  \le
  \sum_y \pi(y)\RewardB(q,y).
\]
\end{theorem}
\begin{proof}
Each term $\pi(y)\RewardB(q,y)$ in the mixture expectation is nonnegative, so
the full finite sum is at least the mass-weighted contribution from
$y_{\mathrm{high}}$:
\[
  \pi(y_{\mathrm{high}})\RewardB(q,y_{\mathrm{high}})
  \le
  \sum_y \pi(y)\RewardB(q,y).
\]
The hypothesis
$\RewardB(q,y_{\mathrm{low}})\le\pi(y_{\mathrm{high}})\RewardB(q,y_{\mathrm{high}})$
then gives the claim by transitivity. In Lean, this is obtained from
\LeanName{B3_weighted_high_lb} (Appendix~\ref{app:b3-scaffolding}).
\end{proof}

\begin{remark}[Why the mass-weighted assumption is necessary]
The stronger informal statement ``positive mass on a higher-entropy feasible
output suffices'' is false in general. A positive probability can be arbitrarily
small, so $\pi(y_{\mathrm{high}})\RewardB(q,y_{\mathrm{high}})$ may still be
less than $\RewardB(q,y_{\mathrm{low}})$. The Lean theorem therefore records
the weakest sufficient algebraic condition used by the proof.
\end{remark}

\begin{theorem}[B3b: Jensen anti-collapse]
\LeanName{B3b_jensen_anti_collapse} via \LeanName{entropy_mixture_concavity}
(PROOF-09, promoted from axiom). For every prompt $q$ and mixture
$\pi:\PMF(Y)$,
\[
  \sum_y \pi(y)H_q(y)
  \le
  \Ent(\mu_q(\pi)).
\]
\end{theorem}
\begin{proof}
Write both sides as sums of $\operatorname{negMulLog}$ terms
(\LeanName{shannonEntropy_eq_sum_negMulLog}). The mixture kernel evaluates
pointwise to the $\pi$-average of component masses
(\LeanName{mixtureKernel_apply_toReal}, from \texttt{PMF.bind\_apply}). For
each downstream state $z$, apply concave Jensen
(\texttt{ConcaveOn.le\_map\_sum} with \texttt{Real.concaveOn\_negMulLog}),
using weights $\pi(y)$ that sum to one (\LeanName{sum_toReal_eq_one}). Summing
over $z$ and interchanging the finite double sum gives the claim.
\end{proof}

\noindent
Consequence: anti-collapse is not only a pairwise dominance property (B1--B2).
At the reward level, if the high-entropy output receives sufficient probability
mass, mixing over feasible outputs cannot do worse than collapsing to a strictly
dominated lower-entropy output. At the distributional level, B3b states that
the entropy of a downstream mixture kernel is at least the $\pi$-weighted average
of component entropies---now proved from primitives, with no axiom dependence.

\subsection{C: Utility Gate, Copy Penalty, and Anti-Padding (Composability Score Only)}

The C1--C3 theorems below remain machine-checked but, after $R_C$ removal,
apply to the Composability Score only; the Benchmark Score contains no
$R_C$ term.

\begin{theorem}[C1: utility gate]
\LeanName{C1_utility_gate}. For any type $\alpha$ with an
\LeanName{AbstractCompressor}, strings $s,\mathcal{H}$, utility value $u$, and
threshold $\tau$, if
\[
  u < \tau,
\]
then
\[
  \RewardC(s,\mathcal{H},u,\tau)=0.
\]
\end{theorem}
\begin{proof}
The definition of $\RewardC$ branches on the condition $\tau\le u$. The
hypothesis $u<\tau$ is equivalent to $\neg(\tau\le u)$, so the false branch of
the definition applies and the reward is $0$. Lean proves this by simplifying
\texttt{rewardC} with \texttt{not\_le.mpr}.
\end{proof}

\noindent
Consequence: novelty alone cannot score; infeasible randomness receives zero
C reward.

\begin{theorem}[C2: copy penalty]
\LeanName{C2_copy_penalty}. Under \LeanName{AbstractCompressor}, for every
corpus $\mathcal{H}$,
\[
  N(\mathcal{H}\mid \mathcal{H}) \le L(\mathcal{H}).
\]
\end{theorem}
\begin{proof}
Unfolding novelty gives
\[
  N(\mathcal{H}\mid\mathcal{H})
  =
  L(\mathcal{H}\mathbin{++}\mathcal{H})-L(\mathcal{H}).
\]
The copy-minimality axiom, applied with $s=\mathcal{H}$ and corpus
$\mathcal{H}$, states exactly
\[
  L(\mathcal{H}\mathbin{++}\mathcal{H})-L(\mathcal{H})
  \le L(\mathcal{H}).
\]
Thus the desired inequality follows directly. Lean unfolds \texttt{novelty} and
uses the copy-minimality field of \LeanName{AbstractCompressor}.
\end{proof}

\noindent
Consequence: exact corpus copies are bounded by the copy-minimality axiom and
cannot be treated as unbounded novelty.

\begin{theorem}[C3: anti-padding]
\LeanName{C3_no_padding_exploit}. Under \LeanName{NormalizedCompressor}, for
any string $s$, padding $p$, and corpus $\mathcal{H}$,
\[
  N(s \mathbin{++} p \mid \mathcal{H})
    \le N(s \mid \mathcal{H}) + L(p).
\]
\end{theorem}
\begin{proof}
Expanding the definition of novelty, the desired inequality becomes
\[
  L(\mathcal{H}\mathbin{++}(s\mathbin{++}p))-L(\mathcal{H})
  \le
  \bigl(L(\mathcal{H}\mathbin{++}s)-L(\mathcal{H})\bigr)+L(p).
\]
After adding $L(\mathcal{H})$ to both sides, this is exactly the normalized
compressor padding-bound axiom:
\[
  L(\mathcal{H}\mathbin{++}(s\mathbin{++}p))
  \le
  L(\mathcal{H}\mathbin{++}s)+L(p).
\]
Lean unfolds \texttt{novelty}, rewrites list associativity, and finishes by
linear arithmetic from the compressor axiom.
\end{proof}

\begin{corollary}[Zero-information padding]
\LeanName{C3_no_trivial_padding_gain}. If additionally $L(p)=0$, then
\[
  N(s \mathbin{++} p \mid \mathcal{H})
    \le N(s \mid \mathcal{H}).
\]
\end{corollary}
\begin{proof}
Apply C3:
\[
  N(s\mathbin{++}p\mid\mathcal{H})
  \le N(s\mid\mathcal{H})+L(p).
\]
Using the additional assumption $L(p)=0$, the right-hand side simplifies to
$N(s\mid\mathcal{H})$. Lean invokes \LeanName{C3_no_padding_exploit} and then
simplifies with $L(p)=0$.
\end{proof}

\noindent
Consequence: the C objective closes the length exploit for padding with zero
information content.

In practice, C is estimated via KL divergence between output and corpus token
PMFs. This is the tractable surface-level approximation of algorithmic code
length; the Lean file proves gating, nonnegativity, and the cross-entropy link
$N_{\mathrm{KL}}+\Ent(p_s)\approx N(s\mid\mathcal{H})$ in
Appendix~\ref{app:n-surface}.

\section{Baseline Contrast: What B and C Add}

The Lean development also proves that B and C improve on natural baselines in
the exact directions claimed.

\subsection{Utility-Only Baseline}

Define a utility-only reward
\[
  R_U(q,y) :=
    \begin{cases}
      U(q,y), & \Feas_m(q,y),\\
      0, & \text{otherwise}.
    \end{cases}
\]

\begin{theorem}[Utility-only indistinguishability]
\LeanName{BBase_indistinguishable}. Assume $Z$ is finite. If $y_1,y_2$ are
feasible, have equal utility,
\[
  U(q,y_1)=U(q,y_2),
\]
and may differ in downstream entropy, then
\[
  R_U(q,y_1)=R_U(q,y_2).
\]
\end{theorem}
\begin{proof}
Since both outputs are feasible, the utility-only reward unfolds to utility on
both sides:
\[
  R_U(q,y_1)=U(q,y_1),\qquad R_U(q,y_2)=U(q,y_2).
\]
The hypothesis $U(q,y_1)=U(q,y_2)$ then gives
$R_U(q,y_1)=R_U(q,y_2)$. The entropy disequality is deliberately unused: it is
the feature the utility-only baseline fails to detect. Lean proves this by
simplifying \texttt{rewardUtil} under feasibility and equal utility.
\end{proof}

\begin{theorem}[B distinguishes what utility-only misses]
\LeanName{BContrast_distinguishes}. Under the same feasibility and equal-utility
assumptions, if
\[
  H_q(y_2) < H_q(y_1),
\]
then
\[
  R_U(q,y_1)=R_U(q,y_2)
  \quad\text{and}\quad
  \RewardB(q,y_2)<\RewardB(q,y_1).
\]
\end{theorem}
\begin{proof}
The first conjunct follows from the previous utility-only argument: feasibility
opens the utility gate and equal utility makes the two baseline rewards equal.
For the second conjunct, feasibility opens the B gate, so
\[
  \RewardB(q,y_i)=H_q(y_i)\qquad(i=1,2).
\]
The entropy hypothesis $H_q(y_2)<H_q(y_1)$ then gives
$\RewardB(q,y_2)<\RewardB(q,y_1)$ by B1. Lean proves the conjunction by
simplifying the first component and applying \LeanName{B1_dominance} for the
second.
\end{proof}

\subsection{Novelty-Only Baseline}

Define an ungated novelty reward
\[
  R_N(s,\mathcal{H}) := N(s\mid\mathcal{H}).
\]

\begin{theorem}[Ungated novelty rewards infeasible novelty]
\LeanName{CBase_infeasible_can_maximize}. Under \LeanName{AbstractCompressor},
if $u<\tau$ and $0<N(s\mid\mathcal{H})$, then
\[
  R_N(s,\mathcal{H})>0
  \quad\text{and}\quad
  \RewardC(s,\mathcal{H},u,\tau)=0.
\]
\end{theorem}
\begin{proof}
By definition, $R_N(s,\mathcal{H})=N(s\mid\mathcal{H})$, so the assumption
$0<N(s\mid\mathcal{H})$ proves $R_N(s,\mathcal{H})>0$. Meanwhile $u<\tau$
means the utility gate in $\RewardC$ is closed, so $\RewardC$ evaluates to $0$.
Lean proves the first conjunct by unfolding \texttt{rewardNovOnly} and the
second by simplifying \texttt{rewardC}.
\end{proof}

\noindent
Consequence: C adds exactly the missing utility gate to novelty-only scoring.

\section{Composability of B, C, and D}

The Lean file \texttt{Creativity/Training/Composability.lean} formalizes the
idea that optimizing B, C, and D jointly creates a Pareto-style objective rather
than a single reducible scalar. The objectives are genuinely multi-dimensional:
Lean contains explicit finite constructions with distinct B, C, and D maximizers
(Appendix~\ref{app:composability-scaffolding}). The main structural fact needed
in the paper is that the Pareto frontier is nonempty.

\begin{theorem}[Nonempty Pareto frontier]
\LeanName{paretoFrontier_nonempty}. For any finite nonempty output type $Y$ and
any rewards $B,C:Y\to\R$, there exists some $y:Y$ that is not jointly dominated
by any other output.
\end{theorem}
\begin{proof}
Choose $y$ to be an argmax of the finite function $B$. Such a maximizer exists
because $Y$ is finite and nonempty. By the previous theorem, no other output can
strictly dominate it in B while being at least as good in C. Therefore this
$y$ lies on the Pareto frontier. Lean obtains the finite maximizer using
\texttt{Finset.exists\_mem\_eq\_sup'}.
\end{proof}


%======================================================================
\part{Benchmark Evaluation Objectives (Frozen, No Gradient Access)}
\label{part:benchmark}
%======================================================================

\noindent
This part collects the benchmark-only quantities: the probe-relative
structural influence $\RewardD$ with its \LeanName{ProbeCompressor} axioms and
full hold-out protocol, the interaction gain $G_k$ with its irreducibility
theorems, the O-information sign bridges, and estimator rank preservation for
B, C, and D. None of these quantities may receive gradient access.

\calloutbox{Goodhart protection}{$\mathcal{P}$ is frozen before all
experimental conditions (Appendix~\ref{app:probe-protocol} hold-out protocol).
$G_k$ is computed post-hoc on fixed agent configurations. Neither
$\mathcal{P}$ nor any benchmark prompt appears in the training data or reward
computation. Violation of this constraint invalidates all rank preservation
guarantees.}

\section{Component D: Probe-Relative Structural Influence}

The file \texttt{Creativity/D/Deformation.lean} extends the compressor with
\LeanName{ProbeCompressor}, adding two axioms:
\begin{itemize}
  \item \textbf{Incorporation helps:} $N(s'\mid\mathcal{H}\cup\{s\})\le
    N(s'\mid\mathcal{H})$ for every probe $s'$;
  \item \textbf{Self-append invariance:} $N(s'\mid\mathcal{H}\mathbin{++}\mathcal{H})
    =N(s'\mid\mathcal{H})$, so exact corpus copies change nothing.
\end{itemize}
The Lean parameter \texttt{probes : List (List $\alpha$)} is a fixed finite
probe set $\mathcal{P}$ satisfying the benchmark hold-out protocol
(Appendix~\ref{app:probe-protocol}). The deformation gain is
\[
  D(s,\mathcal{H},\mathcal{P})
  :=
  \sum_{s'\in\mathcal{P}}
  \bigl(N(s'\mid\mathcal{H})-N(s'\mid\mathcal{H}\cup\{s\})\bigr),
\]
formalized as \LeanName{corpusDeformationGain}. The gated reward is
\LeanName{rewardD}. All D theorems are stated for an arbitrary fixed finite
$\mathcal{P}$; they certify algebraic properties of $D$ given $\mathcal{P}$,
not the choice of $\mathcal{P}$ itself.

\noindent\textbf{Claim framing.}
$D$ measures probe-relative structural influence, not transformational creativity
\emph{by definition}. A random or unrepresentative $\mathcal{P}$ yields
uninformative $D$. When $\mathcal{P}$ is domain-representative---as under the
hold-out protocol below---high $D$ correlates with transformational creativity
in Boden's sense: outputs that compress future domain strings better are
restructuring the domain's information landscape.

\noindent\textbf{Operational probe-set protocol.}
Construct $\mathcal{P}$ as a stratified hold-out sample from the benchmark task
distribution: withhold 10--20\% of the evaluation set before any training run,
then sample $|\mathcal{P}|=100$--$200$ strings from that hold-out stratified by
task category. Freeze $\mathcal{P}$ before all experimental conditions. Record
the sampling seed and split index in Appendix~\ref{app:probe-protocol}.

\begin{theorem}[D1: utility gate]
\LeanName{D1_utility_gate}. If $u<\tau$, then $\RewardD(q,s)=0$.
\end{theorem}
\begin{proof}
Same utility-gating argument as C1 and N1. Lean simplifies \texttt{rewardD}
with \texttt{not\_le.mpr}.
\end{proof}

\begin{theorem}[D2: deformation ordering]
\LeanName{D2_dominance}. If $u\ge\tau$ and
$D(s_1,\mathcal{H},\mathcal{P})<D(s_2,\mathcal{H},\mathcal{P})$, then
$\RewardD(q,s_1)<\RewardD(q,s_2)$.
\end{theorem}
\begin{proof}
Under feasibility, $\RewardD$ equals the ungated deformation gain. The strict
inequality on gains transfers directly. Lean simplifies \texttt{rewardD} with
the feasibility hypothesis and applies the deformation gap.
\end{proof}

\begin{theorem}[D3: exploration bounded by code length (Composability Score only)]
\LeanName{D3_rewardC_bounded}. If $u\ge\tau$, then
\[
  \RewardC(q,s)\le L(s).
\]
\end{theorem}
\begin{proof}
Under feasibility, $\RewardC=N(s\mid\mathcal{H})$. The copy-minimality axiom
of \LeanName{AbstractCompressor} gives exactly
$N(s\mid\mathcal{H})\le L(s)$. Lean unfolds \texttt{rewardC} and applies
\texttt{copy\_minimal}.
\end{proof}

\begin{theorem}[D3: joint exploration--transformation bound (Composability Score only)]
\LeanName{D3_exploration_transformation_bound}. Under
\LeanName{ProbeCompressor}, if $u\ge\tau$, then
\[
  \RewardC(q,s)+\RewardD(q,s)
  \le
  L(s)+D(s,\mathcal{H},\mathcal{P}).
\]
\end{theorem}
\begin{proof}
Under feasibility, the left side is $N(s\mid\mathcal{H})+D(s,\mathcal{H},\mathcal{P})$.
Apply D3\_rewardC\_bounded for the first term and add the nonnegative second
term. Lean rewrites both gated rewards and uses \texttt{linarith}.
\end{proof}

\begin{theorem}[D3 standalone: $R_C$-free deformation bound]
\LeanName{D3_rewardD_bound_standalone}. Under \LeanName{ProbeCompressor},
if $u\ge\tau$, then
\[
  \RewardD(q,s) \le D(s,\mathcal{H},\mathcal{P}).
\]
This is the D3 bound used by the Benchmark Score: no compression-length
($L(s)$/$R_C$) term appears.
\end{theorem}
\begin{proof}
Under feasibility, $\RewardD$ unfolds to the deformation gain itself, so the
bound holds with equality. Lean simplifies \texttt{rewardD} with the
feasibility hypothesis.
\end{proof}

\begin{theorem}[D4: zero probe influence for corpus copies]
\LeanName{D4_zero_for_corpus_copy}. When $s=\mathcal{H}$ is appended to itself,
\[
  D(\mathcal{H},\mathcal{H},\mathcal{P})=0.
\]
\end{theorem}
\begin{proof}
Each probe contribution is
$N(s'\mid\mathcal{H})-N(s'\mid\mathcal{H}\mathbin{++}\mathcal{H})$. The
\LeanName{append_self_invariant} axiom makes every term zero. Lean proves this
by induction on the finite probe list.
\end{proof}

\begin{corollary}[Zero gated reward for corpus copy]
\LeanName{D4_zero_reward_for_corpus_copy}. If $u\ge\tau$ and $s=\mathcal{H}$,
then $\RewardD(q,s)=0$.
\end{corollary}
\begin{proof}
Combine D4 with the feasibility branch of \texttt{rewardD}.
\end{proof}

\noindent
Interpretation: exploratory novelty (C) measures departure from the \emph{current}
corpus. Probe-relative structural influence (D) measures how much the output
\emph{changes} compression of a fixed hold-out probe set $\mathcal{P}$. High C
with $D=0$ is exploratory creativity within fixed rules. High $D$ on a
domain-representative $\mathcal{P}$ correlates with Boden--Wiggins
transformational creativity; the Lean theorems do not equate $D$ with
transformation outright.

\section{Benchmark Estimator Guarantees}

The objective theorems are exact. A benchmark only observes estimates. The
Lean bridge proves that bounded-error estimators preserve rankings when true
gaps exceed twice the error bound.

\subsection{Entropy Estimator}

\begin{definition}[B estimator]
\LeanName{BEstimator}. For a model $m$, a B estimator consists of
\[
  \widehat{H}:Q\to Y\to\R,\qquad \epsilon\in\R,
\]
with uniform error bound
\[
  |\widehat{H}(q,y)-H_q(y)| \le \epsilon
  \quad\text{for all }q,y.
\]
\end{definition}

\begin{theorem}[Entropy rank preservation]
\LeanName{B_rank_preserved}. Assume $Z$ is finite. If $y_1,y_2$ are feasible and
\[
  H_q(y_2)+2\epsilon < H_q(y_1),
\]
then
\[
  \widehat{H}(q,y_2)<\widehat{H}(q,y_1).
\]
\end{theorem}
\begin{proof}
The estimator guarantees give
\[
  |\widehat{H}(q,y_i)-H_q(y_i)|\le\epsilon
  \qquad(i=1,2).
\]
Thus
\[
  \widehat{H}(q,y_2)\le H_q(y_2)+\epsilon
  \quad\text{and}\quad
  H_q(y_1)-\epsilon\le \widehat{H}(q,y_1).
\]
The true-gap assumption
\[
  H_q(y_2)+2\epsilon < H_q(y_1)
\]
rearranges to
\[
  H_q(y_2)+\epsilon < H_q(y_1)-\epsilon.
\]
Chaining these inequalities yields
\[
  \widehat{H}(q,y_2)
  \le H_q(y_2)+\epsilon
  < H_q(y_1)-\epsilon
  \le \widehat{H}(q,y_1).
\]
Lean extracts the one-sided bounds from \texttt{abs\_le.mp} and finishes with
\texttt{linarith}.
\end{proof}

\subsection{Novelty Estimator}

\begin{definition}[C estimator]
\LeanName{CEstimator}. A C estimator consists of
\[
  \widehat{N}:\text{List}(\alpha)\to\text{List}(\alpha)\to\R,
  \qquad \epsilon\in\R,
\]
with uniform error bound
\[
  |\widehat{N}(s,\mathcal{H})-N(s\mid\mathcal{H})| \le \epsilon.
\]
\end{definition}

\begin{theorem}[Novelty rank preservation]
\LeanName{C_rank_preserved}. If
\[
  N(s_2\mid\mathcal{H})+2\epsilon < N(s_1\mid\mathcal{H}),
\]
then
\[
  \widehat{N}(s_2,\mathcal{H})<\widehat{N}(s_1,\mathcal{H}).
\]
\end{theorem}
\begin{proof}
The proof is the novelty analogue of the entropy estimator proof. From the
uniform error bound,
\[
  \widehat{N}(s_2,\mathcal{H})\le N(s_2\mid\mathcal{H})+\epsilon
\]
and
\[
  N(s_1\mid\mathcal{H})-\epsilon
  \le \widehat{N}(s_1,\mathcal{H}).
\]
The true novelty gap
\[
  N(s_2\mid\mathcal{H})+2\epsilon<N(s_1\mid\mathcal{H})
\]
implies
\[
  N(s_2\mid\mathcal{H})+\epsilon
  < N(s_1\mid\mathcal{H})-\epsilon.
\]
Combining these inequalities gives
$\widehat{N}(s_2,\mathcal{H})<\widehat{N}(s_1,\mathcal{H})$. Lean again uses
\texttt{abs\_le.mp} followed by \texttt{linarith}.
\end{proof}

\noindent
Benchmark implication: the harness is reliable for pairwise ranking only when
the true gap exceeds $2\epsilon$. This is an explicit measurement-resolution
condition, not a hidden assumption.


\subsection{Proxy Rank Preservation for D}

The D component deliberately has no trainable estimator of its own: $\RewardD$
is benchmark-only. Its rank-preservation analogue is the conditional proxy
bridge \LeanName{D_tilde_proxy_rank_agreement} (Part~\ref{part:proxy}): under
the \LeanName{ProxyFaithfulness} assumption, a proxy gap
\[
  \RewardDT(s_1)-\RewardDT(s_2)>2\epsilon
\]
forces
$\RewardD(s_1,\mathcal{H},\mathcal{P})>\RewardD(s_2,\mathcal{H},\mathcal{P})$.
Unlike the purely algebraic B and C estimator bounds above, this bridge
carries an \emph{empirical} faithfulness condition that must be monitored at
every training checkpoint (Part~\ref{part:training},
Section~\ref{sec:goodhart-conditions}).
\section{Multi-Agent Irreducibility}

The two-agent Lean model has individual kernels
\[
  T_1(\cdot\mid q,y_1),\qquad T_2(\cdot\mid q,y_2)
\]
and a joint kernel
\[
  T_J(\cdot\mid q,y_1,y_2).
\]
It defines
\[
  H_1=H(T_1),\qquad H_2=H(T_2),\qquad H_J=H(T_J),
\]
and the interaction gain
\[
  G(q,y_1,y_2)
  :=
  H_J(q,y_1,y_2)-\max\{H_1(q,y_1),H_2(q,y_2)\}.
\]
Downstream entropy connects to emergence because it measures whether a
system-level output expands the future state space rather than constraining
subsequent agents into a narrow trajectory. In O-information language,
redundancy repeats the same information across agents, while synergy encodes
information no agent holds alone; B rewards outputs that preserve the possibility
of such recombination.

\begin{assumption}[Faithfulness]
\label{assumption:faithfulness}
Every directed edge $(i,j)$ in the communication graph $\mathcal{G}$ corresponds
to a nonzero pairwise information signal in the joint downstream process: for
any edge $(i,j)$, there exist $q,\mathbf{y}$ such that
\[
  I(z_i;z_j\mid q,\mathbf{y})>0.
\]
This is the standard causal faithfulness assumption
(Spirtes, Glymour, and Scheines, 2000; Pearl, 2009). It rules out degenerate
parameter settings in which agent $j$ structurally depends on agent $i$ but the
statistical effect cancels to zero. Operationally, faithfulness is testable in
the benchmark by measuring pairwise TDMI between agent message streams; TDMI
above the $2\epsilon$ resolution threshold validates the assumption on the
observed task distribution.
\end{assumption}

\begin{theorem}[Corrected zero-gain characterization (PROOF-11)]
\LeanName{gain_zero_iff_joint_eq_max}. For the scalar entropy layer,
\[
  G_k(q,\mathbf{y})=0
  \quad\Longleftrightarrow\quad
  H_J(q,\mathbf{y})=\max_i H_i(q,y_i),
\]
i.e.\ zero gain holds exactly in the degenerate case where the joint process
adds nothing over the best individual agent.
\end{theorem}

\calloutbox{Known gap resolved: the old bridge axiom was false}{The previous
axiom \LeanName{gain_zero_iff_productKernel} asserted
$G_k = 0 \iff \LeanName{isProductKernel}(m)$. This is \emph{incorrect} for
product kernels with positive marginals: entropy additivity
(\LeanName{productKernel_joint_entropy_eq_sum}, now a theorem, PROOF-10) gives
$H_J = H_1 + H_2 > \max(H_1,H_2)$ whenever both marginal entropies are
positive, so the gain is strictly \emph{positive}, not zero. The axiom has
been removed from the Lean development. The refuting direction is
machine-checked as \LeanName{productKernel_additivity_gain_pos}: for a
two-agent coupling with additive joint entropy and positive marginals,
$0 < G_k$. The corrected biconditional \LeanName{gain_biconditional} is
stated against the degeneracy condition $H_J = \max_i H_i$ and is proved, not
assumed.}

\begin{theorem}[Irreducibility from positive interaction gain]
\LeanName{MAS_irreducibility}. Assume $Z$ is finite. If
\[
  0<G(q,y_1,y_2),
\]
then
\[
  H_1(q,y_1)<H_J(q,y_1,y_2)
  \quad\text{and}\quad
  H_2(q,y_2)<H_J(q,y_1,y_2).
\]
\end{theorem}
\begin{proof}
Expanding interaction gain gives
\[
  0 < H_J(q,y_1,y_2)-\max\{H_1(q,y_1),H_2(q,y_2)\}.
\]
Equivalently,
\[
  \max\{H_1(q,y_1),H_2(q,y_2)\}<H_J(q,y_1,y_2).
\]
But each individual entropy is bounded above by the maximum:
\[
  H_1(q,y_1)\le \max\{H_1(q,y_1),H_2(q,y_2)\},
\]
\[
  H_2(q,y_2)\le \max\{H_1(q,y_1),H_2(q,y_2)\}.
\]
Transitivity gives both strict inequalities against $H_J$. Lean unfolds
\texttt{interactionGain} and uses \texttt{le\_max\_left},
\texttt{le\_max\_right}, and \texttt{linarith}.
\end{proof}

\noindent
Consequence: positive interaction gain formally certifies that neither agent
alone achieves the joint downstream entropy. Under
Assumption~\ref{assumption:faithfulness} and the product-kernel bridge below,
positive $G$ reflects genuine multi-agent coordination: zero gain corresponds
to product-kernel independence, while positive gain corresponds to a non-product
joint downstream process.

\begin{definition}[Two-agent estimator]
\LeanName{TwoAgentEstimator}. The MAS benchmark estimator supplies
\[
  \widehat{H}_1,\qquad \widehat{H}_2,\qquad \widehat{H}_J
\]
with a common uniform error bound $\epsilon$:
\[
  |\widehat{H}_1-H_1|\le\epsilon,\quad
  |\widehat{H}_2-H_2|\le\epsilon,\quad
  |\widehat{H}_J-H_J|\le\epsilon.
\]
\end{definition}

\begin{theorem}[Benchmark detection of irreducibility]
\LeanName{MAS_benchmark_detects_irreducibility}. If
\[
  2\epsilon < G(q,y_1,y_2),
\]
then
\[
  \max\{\widehat{H}_1(q,y_1),\widehat{H}_2(q,y_2)\}
  <
  \widehat{H}_J(q,y_1,y_2).
\]
\end{theorem}
\begin{proof}
The estimator bounds imply
\[
  \widehat{H}_1(q,y_1)\le H_1(q,y_1)+\epsilon,\qquad
  \widehat{H}_2(q,y_2)\le H_2(q,y_2)+\epsilon,
\]
and
\[
  H_J(q,y_1,y_2)-\epsilon\le \widehat{H}_J(q,y_1,y_2).
\]
The gain condition
\[
  2\epsilon < H_J(q,y_1,y_2)-\max\{H_1(q,y_1),H_2(q,y_2)\}
\]
rearranges to
\[
  \max\{H_1(q,y_1),H_2(q,y_2)\}+\epsilon
  < H_J(q,y_1,y_2)-\epsilon.
\]
Since $H_1$ and $H_2$ are each at most their maximum, the preceding inequalities
give
\[
  \widehat{H}_1(q,y_1)<\widehat{H}_J(q,y_1,y_2)
  \quad\text{and}\quad
  \widehat{H}_2(q,y_2)<\widehat{H}_J(q,y_1,y_2).
\]
The characterization
\[
  \max\{a,b\}<c \iff a<c \text{ and } b<c
\]
then yields the desired result. Lean extracts the one-sided bounds with
\texttt{abs\_le.mp}, rewrites by \texttt{max\_lt\_iff}, and solves both goals
with \texttt{linarith}.
\end{proof}

\noindent
Benchmark implication: when interaction gain exceeds twice the estimator error,
the harness detects the joint behavior as irreducible.

\subsection{Extension to \texorpdfstring{\(k\)-Agent}{k-Agent} Systems}

The two-agent theorem is the smallest case of a more general Lean development.
For $k>0$ agents, the project defines a \LeanName{KAgentModel} with a
communication graph \LeanName{CommGraph}. A graph edge
\[
  \mathcal{G}(i,j)
\]
means that agent $j$ receives information from agent $i$. The model also records
an abstract pairwise information quantity \LeanName{pairMutualInfo} and a
faithfulness field: every edge has some task/output witness with strictly
positive pairwise information. The individual
kernels
\[
  T_i(\cdot\mid q,y_i)\qquad (i:\operatorname{Fin}(k))
\]
and a joint kernel
\[
  T_J(\cdot\mid q,\mathbf{y}),
  \qquad \mathbf{y}:\operatorname{Fin}(k)\to Y.
\]
The individual and joint entropies are
\[
  H_i(q,y_i) := H(T_i(\cdot\mid q,y_i)),
  \qquad
  H_J(q,\mathbf{y}) := H(T_J(\cdot\mid q,\mathbf{y})).
\]
The $k$-agent interaction gain is
\[
  G_k(q,\mathbf{y})
  :=
  H_J(q,\mathbf{y})
  -
  \max_{i:\operatorname{Fin}(k)} H_i(q,y_i).
\]
This is the direct generalization of the two-agent gain: it compares the joint
downstream entropy to the strongest individual downstream entropy among all
agents.

\begin{theorem}[Base non-product kernel]
\LeanName{nonProduct_base_case}. For two agents, if the communication graph has
an edge $0\to1$, then faithfulness rules out product-kernel independence:
\[
  \neg\,\LeanName{isProductKernel}(m).
\]
\end{theorem}
\begin{proof}
Faithfulness supplies a witness $q,\mathbf{y}$ with positive pairwise
information along the edge. Product-kernel independence forces all off-diagonal
pairwise information values to be zero. These two statements contradict each
other. Lean closes the contradiction by rewriting the product consequence and
using \texttt{linarith}.
\end{proof}

\begin{theorem}[Non-product kernel from an active graph]
\LeanName{connectedMAS_nonProductKernel}. For any number of agents, if
\LeanName{CommGraph} contains at least one edge and the model is faithful, then
\[
  \neg\,\LeanName{isProductKernel}(m).
\]
\end{theorem}
\begin{proof}
Choose one edge $(i,j)$. Irreflexivity gives $i\ne j$. Faithfulness gives
positive pairwise information for that edge at some $q,\mathbf{y}$. A product
kernel forces the same pairwise information value to be zero, contradiction.
No induction on $k$ is needed: one faithful edge is enough.
\end{proof}

\begin{theorem}[Gain biconditional, corrected]
\LeanName{gain_biconditional}. If $0\le G_k(q,\mathbf{y})$, then
\[
  G_k(q,\mathbf{y})=0
  \quad\Longleftrightarrow\quad
  H_J(q,\mathbf{y})=\max_i H_i(q,y_i),
\]
and
\[
  0<G_k(q,\mathbf{y})
  \quad\Longleftrightarrow\quad
  H_J(q,\mathbf{y})\ne\max_i H_i(q,y_i).
\]
\end{theorem}
\begin{proof}
The first equivalence is \LeanName{gain_zero_iff_joint_eq_max}, an algebraic
identity on $G_k = H_J - \max_i H_i$. For the second, nonnegativity leaves
only two cases: $G_k=0$ or $0<G_k$; the zero case is equivalent to the
degeneracy condition, so non-degeneracy is equivalent to strict positivity.
The former product-kernel version of this biconditional relied on the removed
incorrect axiom; see \LeanName{productKernel_additivity_gain_pos} for the
refutation.
\end{proof}

\begin{theorem}[\(k\)-agent irreducibility]
\LeanName{kAgent_irreducibility}. Assume $Z$ is finite and $m$ is a
\LeanName{KAgentModel} with $0<k$. If
\[
  0<G_k(q,\mathbf{y}),
\]
then for every agent $i:\operatorname{Fin}(k)$,
\[
  H_i(q,y_i)<H_J(q,\mathbf{y}).
\]
\end{theorem}
\begin{proof}
Expanding the definition of $G_k$ gives
\[
  0<
  H_J(q,\mathbf{y})
  -
  \max_{i:\operatorname{Fin}(k)}H_i(q,y_i),
\]
so
\[
  \max_{i:\operatorname{Fin}(k)}H_i(q,y_i)
  <
  H_J(q,\mathbf{y}).
\]
For each fixed agent $i$, its individual entropy is bounded by the finite
maximum:
\[
  H_i(q,y_i)\le
  \max_{j:\operatorname{Fin}(k)}H_j(q,y_j).
\]
Transitivity gives
\[
  H_i(q,y_i)<H_J(q,\mathbf{y})
\]
for every $i$. Lean proves this using the finite supremum lemma
\texttt{Finset.le\_sup'} and the definition of
\LeanName{kInteractionGain}.
\end{proof}

\begin{definition}[\(k\)-agent estimator]
\LeanName{KAgentEstimator}. For a fixed model $m$, prompt $q$, and joint action
$\mathbf{y}$, the estimator supplies individual estimates
$\widehat{H}_i$, a joint estimate $\widehat{H}_J$, and a common error bound
$\epsilon$ such that
\[
  |\widehat{H}_i-H_i(q,y_i)|\le\epsilon
  \quad\text{for all }i,
  \qquad
  |\widehat{H}_J-H_J(q,\mathbf{y})|\le\epsilon.
\]
\end{definition}

\begin{theorem}[\(k\)-agent benchmark detection]
\LeanName{kAgent_benchmark_detects}. If
\[
  2\epsilon < G_k(q,\mathbf{y}),
\]
then
\[
  \max_{i:\operatorname{Fin}(k)} \widehat{H}_i
  <
  \widehat{H}_J.
\]
\end{theorem}
\begin{proof}
The proof is the same measurement-resolution argument as in the two-agent case.
For every $i$,
\[
  \widehat{H}_i \le H_i(q,y_i)+\epsilon,
\]
while
\[
  H_J(q,\mathbf{y})-\epsilon\le \widehat{H}_J.
\]
The assumption $2\epsilon<G_k(q,\mathbf{y})$ expands to
\[
  \max_i H_i(q,y_i)+\epsilon
  <
  H_J(q,\mathbf{y})-\epsilon.
\]
Since each $H_i(q,y_i)$ is at most the maximum, each estimate
$\widehat{H}_i$ is strictly below $\widehat{H}_J$. Therefore the finite maximum
of the individual estimates is also strictly below $\widehat{H}_J$. Lean proves
this by rewriting with \texttt{Finset.sup'\_lt\_iff}, extracting estimator
bounds using \texttt{abs\_le.mp}, and closing each agent case with
\texttt{linarith}.
\end{proof}

\noindent
Thus the formal irreducibility criterion is not restricted to two agents. The
two-agent theorem is the readable base case; the $k$-agent theorem is the
general MAS benchmark contract.

\subsection{O-Information Sign Bridges}

The Lean file \texttt{Creativity/MAS/OInformation.lean} also includes scalar
O-information sign checks for the three-agent case. The simplified quantity is
\[
  \Omega_3(H_{12},H_{13},H_{23},H_{123})
  :=
  H_{12}+H_{13}+H_{23}-H_{123}.
\]

\begin{theorem}[Negative O-information from pair-sum bound]
\LeanName{oInformation3_negative_of_pair_sum_lt}. If
\[
  H_{12}+H_{13}+H_{23}<H_{123},
\]
then
\[
  \Omega_3(H_{12},H_{13},H_{23},H_{123})<0.
\]
\end{theorem}
\begin{proof}
Unfolding $\Omega_3$ turns the conclusion into
$H_{12}+H_{13}+H_{23}-H_{123}<0$, which is exactly the assumed strict
inequality rearranged. Lean unfolds \LeanName{oInformation3} and uses
\texttt{linarith}.
\end{proof}

\begin{theorem}[Positive O-information from triple-below-pair-sum]
\LeanName{oInformation3_positive_of_triple_lt_pair_sum}. If
\[
  H_{123}<H_{12}+H_{13}+H_{23},
\]
then
\[
  0<\Omega_3(H_{12},H_{13},H_{23},H_{123}).
\]
\end{theorem}
\begin{proof}
After unfolding $\Omega_3$, the conclusion is
$0<H_{12}+H_{13}+H_{23}-H_{123}$, which follows by rearranging the hypothesis.
Lean again closes the arithmetic by \texttt{linarith}.
\end{proof}

\begin{theorem}[Equal-marginals synergy bridge]
\LeanName{equalMarginals_gain_implies_synergy}. If $h<H_J$ and the additional
sign condition
\[
  3(H_J-h)<H_J
\]
holds, then
\[
  \Omega_3(H_J-h,H_J-h,H_J-h,H_J)<0.
\]
\end{theorem}
\begin{proof}
Substituting the equal-marginals expression into $\Omega_3$ gives
\[
  3(H_J-h)-H_J.
\]
The sign condition says this quantity is negative. Lean unfolds
\LeanName{oInformation3} and uses \texttt{linarith}. The explicit sign
condition is necessary: $h<H_J$ alone does not force $3(H_J-h)<H_J$.
\end{proof}

\begin{theorem}[Negative-gain redundancy bridge]
\LeanName{negativeGain_implies_redundancy}. If $H_J<h$ and
\[
  H_J<3(h-H_J),
\]
then
\[
  0<\Omega_3(h-H_J,h-H_J,h-H_J,H_J).
\]
\end{theorem}
\begin{proof}
Unfolding $\Omega_3$ reduces the conclusion to
$0<3(h-H_J)-H_J$, which is exactly the second hypothesis rearranged. Lean
finishes by \texttt{linarith}.
\end{proof}

\subsection{O-Information Kernel Bridge}

The file \texttt{Creativity/MAS/OInfoBridge.lean} connects the scalar
O-information sign lemmas to concrete product-kernel structure. A
\LeanName{TwoAgentProductModel} specifies independent kernels $k_1,k_2$ and
forms their joint PMF by independent coupling.

\begin{theorem}[Product-kernel entropy additivity (PROOF-10, promoted from axiom)]
\LeanName{productKernel_joint_entropy_eq_sum}. For finite downstream alphabet
$Z$ and independent product coupling,
\[
  \Ent(Z_1,Z_2)
  =
  \Ent(Z_1)+\Ent(Z_2),
\]
where $(Z_1,Z_2)$ are drawn from the product kernel induced by $k_1(q,y_1)$ and
$k_2(q,y_2)$.
\end{theorem}
\begin{proof}
Formerly an axiom; now proved as \LeanName{shannonEntropy_jointPMF}. The
product coupling factorizes pointwise, $(p\otimes q)(a,b)=p(a)q(b)$
(\LeanName{jointPMF_apply}, via \texttt{PMF.bind\_apply} and
\texttt{PMF.map\_apply}). Expanding the joint entropy over $Z\times Z$ and
applying the logarithm product rule in the form
$\operatorname{negMulLog}(xy)=y\operatorname{negMulLog}(x)+x\operatorname{negMulLog}(y)$
(\texttt{Real.negMulLog\_mul}) splits the double sum into
$\sum_a\operatorname{negMulLog}(p(a))\sum_b q(b)
 +\sum_b\operatorname{negMulLog}(q(b))\sum_a p(a)$;
the mass identities $\sum p = \sum q = 1$ (\LeanName{sum_toReal_eq_one})
reduce this to $\Ent(Z_1)+\Ent(Z_2)$.
\end{proof}

\begin{theorem}[Product kernels yield positive pair-sum entropy]
\LeanName{productKernel_implies_redundancy}. If
$0<\Ent(k_1(q,y_1))$ and $0<\Ent(k_2(q,y_2))$, then
\[
  0<\Ent(k_1(q,y_1))+\Ent(k_2(q,y_2)).
\]
\end{theorem}
\begin{proof}
This is pure arithmetic: the sum of two strictly positive reals is strictly
positive. Lean closes with \texttt{linarith}.
\end{proof}

\begin{theorem}[Product-kernel gain is not synergy]
\LeanName{productKernel_gain_is_not_synergy}. Independent product kernels can
have positive individual and joint entropies without satisfying a negative
O-information condition. The Lean development records this as a scoped boundary
theorem: product independence is not the same object as irreducible synergy.
\end{theorem}
\begin{proof}
Lean proves this documentation theorem by \texttt{trivial}. The mathematical
content is the conceptual separation between product-kernel additivity and the
sign conditions required for negative O-information in the scalar bridge
lemmas above.
\end{proof}

\section{Interaction Gain to Decision Quality: the CUE Bridge}

The entropy-level irreducibility theorems above certify that positive $G_k$
means no single agent reproduces the joint downstream process. The following
bridge (PROOF-07) lifts this to the receiver decision layer of
Part~\ref{part:cue}.

\begin{definition}[Receiver calibration link]
\LeanName{ReceiverCalibratedMAS}. A receiver whose CUE readings are calibrated
to downstream entropy with uniform error $\varepsilon_{\mathrm{cal}}$: for
each agent $i$, $|\mathrm{CUE}_i - H_i| \le \varepsilon_{\mathrm{cal}}$, and
for the joint output, $|\mathrm{CUE}_J - H_J| \le \varepsilon_{\mathrm{cal}}$.
This is the formal \texttt{brier\_entropy\_calibration\_link}, parallel in
structure to \LeanName{KAgentEstimator}.
\end{definition}

\begin{theorem}[PROOF-07: interaction gain forces CUE improvement]
\LeanName{Gk_implies_CUE_improvement} and
\LeanName{Gk_implies_CUE_improvement_sup}. Under the receiver calibration
link, if
\[
  2\varepsilon_{\mathrm{cal}} < G_k(q,\mathbf{y}),
\]
then $\mathrm{CUE}_i < \mathrm{CUE}_J$ for every agent $i$, and hence
\[
  \max_i \mathrm{CUE}(y_i, A, T) < \mathrm{CUE}(y_{\mathrm{joint}}, A, T).
\]
\end{theorem}
\begin{proof}
The CUE analogue of \LeanName{kAgent_benchmark_detects}: each agent CUE is at
most $H_i + \varepsilon_{\mathrm{cal}} \le \max_i H_i +
\varepsilon_{\mathrm{cal}}$, the joint CUE is at least
$H_J - \varepsilon_{\mathrm{cal}}$, and the gain hypothesis separates the two
bounds. Lean closes with \texttt{linarith} over the absolute-value bounds.
\end{proof}

\section{Trajectory-Level Diagnostics: Step-CUE and Diverge-Converge}

These benchmark diagnostics evaluate the \emph{reasoning trajectory}
$y_{1:t}$ rather than only the final output.

\begin{theorem}[PROOF-04: gated Step-CUE curves are monotone]
\LeanName{CUE_trajectory_monotone_nondecreasing}. Let
$C(t) = \mathrm{CUE}(y_{1:t}, A, T)$ be the running Step-CUE curve and
$r_t = C(t+1) - C(t)$ the incremental step reward (\LeanName{stepReward}). If
every step satisfies the utility gate ($r_t \ge 0$ for all $t$), then $C$ is
non-decreasing.
\end{theorem}
\begin{proof}
Nonnegative successive differences give monotonicity;
\texttt{monotone\_nat\_of\_le\_succ} plus \texttt{linarith}.
\end{proof}

\begin{theorem}[PROOF-14: $\gamma$-positivity characterization]
\LeanName{step_CUE_gamma_positivity_characterization}. For the fitted
Step-CUE curve $C(t) = C_\infty(1 - e^{-\mu t}) + \gamma t$
(\LeanName{stepCUEFit}), the linear saturation rate satisfies
\[
  \gamma > 0
  \iff
  \exists\, t > 0:\;
  C_\infty(1 - e^{-\mu t}) < C(t),
\]
i.e.\ $\gamma$ is positive exactly when some step exceeds the pure
exponential-saturation prediction --- genuine continuous discovery.
\end{theorem}
\begin{proof}
The fitted curve exceeds its exponential component by exactly $\gamma t$;
positivity of $\gamma t$ at a positive time is equivalent to positivity of
$\gamma$. Lean closes both directions with \texttt{linarith}/\texttt{nlinarith}.
\end{proof}

\begin{theorem}[PROOF-05: Diverge-Converge score existence and non-triviality]
\LeanName{DC_score_existence} and \LeanName{monotone_increasing_not_DC}.
There exists a receiver-entropy trajectory $E$ and a peak index $t^\ast$
satisfying the diverge-converge predicate
(\LeanName{hasDivergeConverge}): $E$ strictly increases up to $t^\ast$
(divergence) and strictly decreases afterward (convergence), so
$\mathrm{DC} = 1$ is achievable. Conversely, any strictly increasing
(monotone-divergent) trajectory fails the predicate for every choice of
$t^\ast$, so $\mathrm{DC} = 0$ on monotone trajectories: the score is neither
trivially satisfied nor trivially failed.
\end{theorem}
\begin{proof}
Existence is by explicit construction: the tent-shaped trajectory
$E(t) = 1 - (t-1)^2$ on three steps peaks at $t^\ast = 1$
(\LeanName{tentTrajectory}). Non-triviality: the convergence phase at the peak
would contradict monotonicity at $t^\ast$; \texttt{linarith}.
\end{proof}

\section{Estimator Calibration and Benchmark Specification}

The Lean benchmark layer contains one imported statistical axiom and one
machine-readable experiment specification. The plug-in entropy estimator is
defined directly from empirical counts:
\[
  \widehat{H}_{\mathrm{plug}}(c,n)
  =
  -\sum_y \frac{c_y}{n}\log\frac{c_y}{n}.
\]

\begin{assumption}[Plug-in entropy bias bound (PROOF-12; Known Gaps Registry KG-2)]
\LeanName{plugIn_bias_axiom}. For alphabet $Y$, counts $c:Y\to\mathbb{N}$,
sample size $n>0$, and true entropy $H$, the imported Miller--Madow-scale bound
is
\[
  \left|\widehat{H}_{\mathrm{plug}}(c,n)-H\right|
  \le
  \frac{|Y|}{2n}.
\]
This is a \emph{registered trusted import} (Miller \& Madow, 1955), not an
internally proved theorem. Proof target for future promotion: Hoeffding's
inequality applied to the empirical frequency estimator, combined with the
Lipschitz continuity of $x\mapsto -x\log x$ on $[0,1]$ (each term bounded in
$[0,1/e]$); a Lean~4 statistical-learning-theory library providing
\texttt{Hoeffding\_inequality} would discharge it. The subsequent Lean
corollary uses the bound algebraically.
\end{assumption}

\begin{theorem}[Sample requirement for entropy ranking]
\LeanName{entropy_ranking_sample_requirement}. If
\[
  H_{\mathrm{low}}+\frac{|Y|}{n}<H_{\mathrm{high}},
\]
then the plug-in estimates preserve the ordering:
\[
  \widehat{H}_{\mathrm{plug}}(c_{\mathrm{low}},n)
  <
  \widehat{H}_{\mathrm{plug}}(c_{\mathrm{high}},n).
\]
\end{theorem}
\begin{proof}
The bias axiom gives each estimate error at most $|Y|/(2n)$. Therefore the
total possible two-sided error between the two estimates is at most $|Y|/n$.
If the true gap exceeds this amount, the lower estimate remains below the higher
estimate even in the worst case. Lean extracts the absolute-value bounds with
\texttt{abs\_le.mp}, proves
\[
  2\cdot \frac{|Y|}{2n}=\frac{|Y|}{n},
\]
and finishes by \texttt{linarith}.
\end{proof}

To make the resolution condition concrete: for a 50-token effective alphabet and
$n=200$ episodes, the calibration gives
\[
  \epsilon=\frac{|Y|}{2n}=\frac{50}{400}=0.125
\]
nats, so the benchmark's two-sided detection floor is
$2\epsilon=0.25$ nats. Gaps below $0.25$ nats require larger samples; for
example, reaching a $0.10$-nat floor with the same effective alphabet requires
$n\ge 500$ episodes. All reported experiments should state their effective
alphabet size, sample count, and minimum observed entropy gap against this
calibration floor.

The experiment specification file defines four topology labels
\LeanName{Topology}, an episode record \LeanName{EpisodeResult}, and two
predicates:
\LeanName{predictedS2DominatesS1} and
\LeanName{predictedGroupthinkSignature}.

\begin{theorem}[Detectability of S2 dominance]
\LeanName{benchmark_can_detect_S2_dominance}. If the true interaction gain gap
between decentralized S2 and dense-centralized S1 exceeds two resolution margins,
\[
  G_{\mathrm{S1}}+4\epsilon<G_{\mathrm{S2}},
\]
then there exist benchmark episode records $r_2,r_1$ with topologies S2 and S1
respectively such that
\[
  \text{\LeanName{predictedS2DominatesS1}}(r_2,r_1)
\]
holds.
\end{theorem}
\begin{proof}
Lean constructs worst-case estimated gains:
\[
  \widehat{G}_{\mathrm{S2}}=G_{\mathrm{S2}}-\epsilon,
  \qquad
  \widehat{G}_{\mathrm{S1}}=G_{\mathrm{S1}}+\epsilon.
\]
The predicate requires
\[
  \widehat{G}_{\mathrm{S1}}+2\epsilon<\widehat{G}_{\mathrm{S2}},
\]
which becomes
\[
  G_{\mathrm{S1}}+3\epsilon<G_{\mathrm{S2}}-\epsilon,
\]
equivalently $G_{\mathrm{S1}}+4\epsilon<G_{\mathrm{S2}}$. This is exactly the
hypothesis. Lean fills the existential records explicitly and closes the
arithmetic with \texttt{linarith}.
\end{proof}

\paragraph{Benchmark interpretation.}
The benchmark operationalizes the formal results by checking whether systems
optimizing B maintain richer future trajectory spaces, whether systems optimizing
C reduce corpus-adjacent repetition without rewarding incoherence or padding,
and whether multi-agent systems with detected interaction gain behave
irreducibly rather than as sums of individual agents. The theory says what
should happen; the benchmark tests whether it happens under the stated
resolution and faithfulness assumptions.


%======================================================================
\part{Dual-Process Separation Principle}
\label{part:dps}
%======================================================================

\noindent
This part sits deliberately \emph{between} the frozen benchmark metrics
(Part~\ref{part:benchmark}) and the training signal
(Part~\ref{part:training}): the separation proved here is what makes the
benchmark non-Goodhartable. Goodhart resistance in this framework is a
\textbf{two-level structure}:

\begin{itemize}
  \item \textbf{Level 1 --- Architectural Goodhart resistance.} The
  Dual-Process Separation Principle (PROOF-03): the training gradient is
  structurally independent of every frozen benchmark object.
  \item \textbf{Level 2 --- Protocol Goodhart resistance.} The four protocol
  conditions collected in \LeanName{Goodhart_resistance_summary}
  (Part~\ref{part:training}), each paired with its collapse direction by the
  Structural Goodhart Theorem (PROOF-13) below.
\end{itemize}

\section{The Separation Theorem}

Let $\Pi_{\mathrm{frozen}} = \{\mathcal{P}, T, A\}$ be the frozen benchmark
objects (probe set, task battery, receiver) and $\Pi_{\mathrm{live}} =
\{\RewardDT\text{-probes}, \text{policy}, \text{corpus schedule}\}$ the live
training quantities. A training configuration splits into a live and a frozen
component (\LeanName{TrainingConfig}); a reward \emph{depends only on the live
component} when any two configurations agreeing on the live part receive equal
reward (\LeanName{dependsOnlyOnLive}) --- the discrete statement that its
gradient in the frozen coordinate is zero.

\begin{theorem}[PROOF-03: Dual-Process Separation Principle]
\LeanName{dual_process_separation_principle}. Any training reward built as a
function of the live component alone (\LeanName{liveReward}) depends only on
the live component: for every gradient update step,
$\nabla_\theta \RewardTrain \perp \Pi_{\mathrm{frozen}}$.
\end{theorem}

\begin{corollary}[No differentiable Goodhart path]
\LeanName{Rtrain_cannot_goodhart_frozen}. Changing any frozen benchmark object
while holding the live component fixed leaves $\RewardTrain$ unchanged.
Consequently no gradient step on $\RewardTrain$ can be driven by, or exploit,
a frozen object: $\RewardTrain$ cannot Goodhart the benchmark through any
differentiable path. Combined with \LeanName{trainReward_neq_fullReward}, the
train/benchmark split is structural, not merely definitional.
\end{corollary}

\section{The Structural Goodhart Collapse Theorem}

Protection and violation are now a theorem pair: the rank-agreement guarantee
(\LeanName{D_tilde_proxy_rank_agreement}, Part~\ref{part:proxy}) holds under
the protocol conditions, and this section proves each violation admits a
refuting exploitation path.

\begin{definition}[Rank agreement and rank reversal]
\LeanName{rankAgreement}, \LeanName{rankReversal}. For a proxy/benchmark
instance with resolution margin $\varepsilon$, rank agreement demands that a
proxy gap beyond $2\varepsilon$ forces the same benchmark ordering; a rank
reversal is a pair $s_1, s_2$ where the proxy prefers $s_1$ beyond
$2\varepsilon$ while the benchmark strictly prefers $s_2$.
\end{definition}

\begin{theorem}[PROOF-13: Structural Goodhart Collapse]
\LeanName{rank_reversal_refutes_agreement},
\LeanName{reversal_instance_exists},
\LeanName{structural_goodhart_collapse}. Any rank reversal refutes rank
agreement, and rank reversals are realizable: there is an explicit
two-candidate instance whose proxy ranks the candidates oppositely to the
benchmark beyond the margin. Hence whenever a violated Goodhart condition
admits a reversal, the proxy-rank-agreement guarantee collapses.
\end{theorem}

\begin{remark}[The four exploitation paths]
Each protocol condition instantiates the collapse through its own path:
(i) \emph{fixed probe target} --- if $Q_{\mathrm{probe}}$ contains a fixed
string across steps, the policy can memorize it, inflating $\RewardDT$ without
improving $R_D$ on the frozen $\mathcal{P}$;
(ii) \emph{corpus non-stationarity} --- if $\mathcal{H}$ absorbs policy
outputs, the anti-padding baseline shifts and \LeanName{C3_no_padding_exploit}
loses its precondition;
(iii) \emph{loose utility gate} --- if $\tau \le H_q(y_{\mathrm{random}})$,
incoherent high-entropy outputs pass the gate
(\LeanName{CBase_infeasible_can_maximize} pattern);
(iv) \emph{undetected proxy decay} --- if $\rho(\RewardDT, R_D) < \rho_{\min}$
without detection, the faithfulness premise of rank agreement fails. In each
case the admitted reversal triggers
\LeanName{rank_reversal_refutes_agreement}.
\end{remark}

%======================================================================
\part{2-CGRPO Training Signal (Certified, Goodhart-Resistant)}
\label{part:training}
%======================================================================

\noindent
This part defines the reward the trainer actually optimizes, proves that DPO
steps ascend each of its components, and states the four Goodhart-resistance
conditions that a training run must satisfy for the certificates to apply.

\section{Training Reward Definition}
\label{sec:train-reward}

The training reward is
\[
  \RewardTrain(q,y,s)
  :=
  \RewardB(q,y)
  +\lambda_{\tilde D}\,\RewardDT(q,s)
  \qquad (\text{\LeanName{trainReward}}),
\]
which is a distinct object from the reported multiplicative Benchmark Score
$R_{\mathrm{creativity}}$~\eqref{eq:Rcreativity}
(\LeanName{Rcreativity}, Section~\ref{sec:benchmark-score}).
There is no $\lambda_C\RewardC$ channel. The components are:
\begin{itemize}
  \item $\RewardB$: utility-gated downstream entropy, exactly as in
    Part~\ref{part:definition}.
  \item $\RewardDT$: the trainable proxy for D, defined below on
    \emph{resampled} probe prompts, never on the frozen benchmark probe set
    $\mathcal{P}$.
\end{itemize}
The train/benchmark split is machine-checked by
\LeanName{trainReward_neq_benchmarkReward}.

\subsection{The Trainable Proxy \texorpdfstring{$\RewardDT$}{R-tilde-D}}

The Lean structure \LeanName{DTildeModel}
(\texttt{Creativity/D/ProxyReward.lean}) contains a probe prompt distribution
$Q_{\mathrm{probe}}:\PMF(Q)$ that is resampled per batch and never frozen, a
frozen reference policy $p_{\mathrm{ref}}:Q\to\PMF(Z)$ that is never updated
during training, and the trained policy's conditioned predictive distribution
$p(\cdot\mid q_{\mathrm{probe}},s)$. The proxy reward is
\[
  \RewardDT(s,q_{\mathrm{probe}})
  =
  D_{\mathrm{KL}}\!\left(
    p(\cdot \mid q_{\mathrm{probe}}, s)
    \;\Big\|\;
    p_{\mathrm{ref}}(\cdot \mid q_{\mathrm{probe}})
  \right)
  \cdot \mathbf{1}[u \geq \tau],
\]
formalized as \LeanName{rewardDTilde}.

\begin{theorem}[$\tilde D$ utility gate]
\LeanName{D_tilde_utility_gate}. If $u<\tau$, then
$\RewardDT(s,q_{\mathrm{probe}})=0$.
\end{theorem}
\begin{proof}
Mirrors \LeanName{D1_utility_gate}: the gate closes the reward branch. Lean
simplifies \texttt{rewardDTilde} with \texttt{not\_le.mpr}.
\end{proof}

\begin{theorem}[$\tilde D$ nonnegativity]
\LeanName{D_tilde_nonneg}. For all $s,q_{\mathrm{probe}},u,\tau$,
$0\le\RewardDT(s,q_{\mathrm{probe}})$.
\end{theorem}
\begin{proof}
The gated branch is a KL divergence, nonnegative by
\LeanName{klDivergence_nonneg}; the closed branch is $0$.
\end{proof}

\begin{theorem}[Zero proxy reward for reference-identical policies]
\LeanName{D_tilde_zero_for_identical}. If
$p(\cdot\mid q_{\mathrm{probe}},s)=p_{\mathrm{ref}}(\cdot\mid q_{\mathrm{probe}})$,
then $\RewardDT(s,q_{\mathrm{probe}})=0$.
\end{theorem}
\begin{proof}
Mirrors \LeanName{D4_zero_for_corpus_copy}: reproducing the reference
distribution earns nothing. Lean rewrites with the identity hypothesis and
the term-by-term lemma \LeanName{klDivergence_self}.
\end{proof}

\begin{assumption}[No fixed probe target]
\LeanName{D_tilde_no_fixed_target}. $Q_{\mathrm{probe}}$ is drawn fresh each
step: no string in $Q_{\mathrm{probe}}$ is fixed across training steps.
Formally, for the probe batch schedule of any admissible protocol
(\LeanName{DTildeTrainingProtocol}), no probe prompt belongs to the batch at
every step. This is the formal Goodhart-resistance condition for
$\RewardDT$: it prevents memorization of specific probe inputs.
\end{assumption}

\begin{theorem}[$\tilde D$ preference ordering]
\LeanName{D_tilde_dominance}. If $u\ge\tau$ and
$\RewardDT(s_1)<\RewardDT(s_2)$, then $(s_2,s_1)$ is a $\tilde D$-valid
preference pair.
\end{theorem}
\begin{proof}
Same structure as \LeanName{B1_dominance} and \LeanName{D2_dominance}: under
the open gate the reward unfolds to the KL shift, so the strict ordering
transfers. This is the prerequisite for
\LeanName{DPO_D_tilde_gradient_correct} below.
\end{proof}

\subsection{Properties of \texorpdfstring{$\RewardTrain$}{R\_train}}

\begin{theorem}[Training reward nonnegativity]
\LeanName{trainReward_nonneg}. If $0\le\RewardB$, $0\le\RewardC$,
$0\le\lambda_C$, and $0\le\lambda_{\tilde D}$, then
$0\le\RewardTrain(q,y,s)$.
\end{theorem}
\begin{proof}
Mirrors \LeanName{benchmarkReward_nonneg}; the $\tilde D$ component is
nonnegative unconditionally by \LeanName{D_tilde_nonneg}.
\end{proof}

\begin{theorem}[Training advantage sign correctness]
\LeanName{trainReward_advantage_sign_correct}. For a group of $\RewardTrain$
scores, if $r_j<\mu<r_i$ and the stabilizer $\epsilon>0$, then $A_j<0$ and
$0<A_i$.
\end{theorem}
\begin{proof}
Mirrors \LeanName{advantage_sign_correct}, instantiated at the training
reward via \LeanName{trainRewards}.
\end{proof}

\begin{theorem}[Train/benchmark split]
\LeanName{trainReward_neq_fullReward}. $\RewardTrain\ne\RewardFull$: whenever
$\lambda_{\tilde D}\RewardDT(s)\ne\lambda_D\RewardD(s,\mathcal{H},\mathcal{P})$,
the two rewards differ at that point. $\RewardD$ and $\RewardDT$ are distinct
objects; this theorem exists solely to make the train/benchmark split
machine-checkable, not just a prose convention.
\end{theorem}
\begin{proof}
Both rewards share the B and C terms, so their difference is exactly the
difference of the weighted D components. Lean unfolds both definitions and
closes with \texttt{linarith}.
\end{proof}

\section{Training Correspondence: DPO Direction}

The Lean development proves a local DPO-B training guarantee at the abstraction
level used in the finite model. The calculus-heavy softmax-gradient step is
represented by two explicit local hypotheses: \textbf{H1} non-degeneracy, so the
winner and loser log-probability gradients are not identical, and \textbf{H2} a
sufficiently small step size, so the local DPO step transfers positive
probability mass from the B-loser to the B-winner while leaving other outputs
fixed to first order. Under this local transfer hypothesis, Lean proves that
expected B reward strictly increases.

For reference, the older bridge file \texttt{Creativity/B/DPO.lean} also defines
a finite softmax policy \LeanName{softmaxPolicy} and DPO log-odds
$\operatorname{LogOdds}_\ell(y_w,y_l)=\ell_{y_w}-\ell_{y_l}$.

\begin{theorem}[Finite logit update increases winner log-odds]
\LeanName{dpo_logit_direction}. Assume $Y$ has decidable equality. Let
$\ell:Y\to\R$, distinct winner and loser indices $y_w\ne y_l$, and step
$\eta>0$. Define the updated logits
\[
  \ell'_y =
  \begin{cases}
    \ell_y+\eta, & y=y_w,\\
    \ell_y-\eta, & y=y_l,\\
    \ell_y, & \text{otherwise}.
  \end{cases}
\]
Then
\[
  \operatorname{LogOdds}_\ell(y_w,y_l)
  <
  \operatorname{LogOdds}_{\ell'}(y_w,y_l).
\]
\end{theorem}
\begin{proof}
Because $y_w\ne y_l$, the update gives $\ell'_{y_w}=\ell_{y_w}+\eta$ and
$\ell'_{y_l}=\ell_{y_l}-\eta$. Hence
\[
  \operatorname{LogOdds}_{\ell'}(y_w,y_l)
  =
  (\ell_{y_w}+\eta)-(\ell_{y_l}-\eta)
  =
  \operatorname{LogOdds}_\ell(y_w,y_l)+2\eta.
\]
Since $\eta>0$, the strict increase follows. Lean packages the update in
\LeanName{LogitUpdate}/\LeanName{applyUpdate}, simplifies with the disequality
hypothesis, and closes with \texttt{linarith}.
\end{proof}

\begin{theorem}[DPO gradient alignment]
\LeanName{dpo_gradient_increases_logodds}. If $\beta>0$,
$\sigma(r_\theta(y_l)-r_\theta(y_w))>0$, the squared gradient-difference norm is
positive, and $\eta>0$, then the first-order DPO update has strictly positive
alignment with the winner-over-loser log-odds direction.
\end{theorem}
\begin{proof}
Lean formalizes the alignment scalar as
\[
  \beta\,\sigma(r_l-r_w)\,
  \|\nabla\log\pi_\theta(y_w)-\nabla\log\pi_\theta(y_l)\|^2.
\]
All factors are strictly positive under H1 and $\beta>0$, so multiplying by
$\eta>0$ remains strictly positive. Lean closes this with positivity.
\end{proof}

\begin{theorem}[Softmax monotonicity for the winning logit]
\LeanName{softmax_monotone_in_logit_diff}. For two competing logits, increasing
the winner logit by any $\delta>0$ strictly increases the winner's softmax
probability.
\end{theorem}
\begin{proof}
Writing $a=\exp(\ell_w)>0$, $b=\exp(\ell_l)>0$, and $c=\exp(\delta)>1$, the
claim is
\[
  \frac{a}{a+b}<\frac{ac}{ac+b}.
\]
Cross-multiplication over positive denominators reduces this to
$ab<acb$, which follows from $c>1$. Lean proves the inequality by exponential
positivity, cross-multiplication, and ring normalization.
\end{proof}

\begin{theorem}[B-valid probability transfer increases expected B]
\LeanName{dpo_step_increases_expected_RB}. Let a local DPO step transfer
$\delta>0$ probability mass from $y_l$ to $y_w$, holding all other output
probabilities fixed. If
\[
  R_B(q,y_l)<R_B(q,y_w),
\]
then the policy expectation of $R_B$ strictly increases.
\end{theorem}
\begin{proof}
The expectation change is exactly
\[
  \delta R_B(q,y_w)-\delta R_B(q,y_l)
  =
  \delta\bigl(R_B(q,y_w)-R_B(q,y_l)\bigr).
\]
This is positive because $\delta>0$ and the pair is B-valid. Lean proves the
finite-sum identity in \LeanName{expectedReward_pairwiseTransfer_eq} and then
uses \texttt{linarith}.
\end{proof}

\begin{theorem}[DPO-B gradient correctness]
\LeanName{DPO_B_gradient_correct}. Under H1 and H2, encoded as the local DPO
probability-transfer hypothesis, a single DPO step on a B-valid preference pair
strictly increases expected downstream entropy:
\[
  \mathbb{E}_{y\sim\pi_{\theta'}}[R_B(q,y)]
  >
  \mathbb{E}_{y\sim\pi_{\theta}}[R_B(q,y)].
\]
\end{theorem}
\begin{proof}
The local DPO hypothesis supplies a $\delta>0$ transfer from the B-loser to the
B-winner. The preceding theorem converts this transfer and B-validity directly
into a strict increase in expected $R_B$. This is the full local training
correspondence proved in Lean; the finite logit theorem above is the two-logit
special case.
\end{proof}

\begin{theorem}[B-valid DPO pair moves in the correct direction]
\LeanName{DPO_B_correct_direction}. Assume $Z$ is finite and $Y$ has decidable
equality. Let $y_w,y_l:Y$ be feasible and satisfy
\[
  H_q(y_l)<H_q(y_w).
\]
For any logits $\ell:Y\to\R$ and step size $\eta>0$, define $\ell'$ as above.
Then
\[
  \operatorname{LogOdds}_\ell(y_w,y_l)
  <
  \operatorname{LogOdds}_{\ell'}(y_w,y_l).
\]
\end{theorem}
\begin{proof}
The entropy dominance hypothesis $H_q(y_l)<H_q(y_w)$ implies $y_w\ne y_l$;
otherwise it would reduce to $H_q(y_w)<H_q(y_w)$, impossible by irreflexivity
of $<$. The remainder is exactly \LeanName{dpo_logit_direction} applied to the
same finite update. Lean derives the disequality from the entropy gap, simplifies
the if-then-else update, and closes with \texttt{linarith}.
\end{proof}

\noindent
Training-bridge summary: B-constructed preference pairs increase winner
log-odds under the finite two-logit update, and under the explicit local
non-degeneracy/small-step transfer hypotheses, a DPO step strictly increases
expected $R_B$. This justifies the DPO-style training language as a local ascent
claim for the B objective.

\section{RL Reward Signal Bridge}

The Lean file \texttt{Creativity/Training/RLSignal.lean} connects the abstract
B and D rewards to a group-relative RL signal. It defines the additive B--D
diagnostic
\[
  R_{BD}(q,y,s)
  =
  \RewardB(q,y)+\lambda_D\RewardD(s,\mathcal{H},u,\tau)
  \qquad (\text{\LeanName{benchmarkReward}}),
\]
which serves the B/D rank-preservation correspondence. The \emph{reported}
Benchmark Score is the multiplicative \LeanName{Rcreativity} of
Section~\ref{sec:benchmark-score}~\eqref{eq:Rcreativity}. For a finite group
of rewards $r_i$, the group-relative advantage is
\[
  A_i
  =
  \frac{r_i-\mu}{\sqrt{\frac{1}{G}\sum_j(r_j-\mu)^2}+\epsilon},
  \qquad
  \mu=\frac{1}{G}\sum_j r_j.
\]

\begin{theorem}[B--D diagnostic nonnegativity]
\LeanName{benchmarkReward_nonneg}. If $0\le\RewardB$, $0\le\RewardD$, and
$0\le\lambda_D$, then $0\le R_{BD}(q,y,s)$.
\end{theorem}
\begin{proof}
Sum of a nonnegative term and a product of nonnegatives
(\texttt{add\_nonneg}, \texttt{mul\_nonneg}).
\end{proof}

\begin{theorem}[Positive RL denominator]
\LeanName{groupDenom_pos}. If $\epsilon>0$, then
\[
  \sqrt{\frac{1}{G}\sum_j(r_j-\mu)^2}+\epsilon>0.
\]
\end{theorem}
\begin{proof}
The square-root term is always nonnegative. Adding the strictly positive
$\epsilon$ makes the denominator strictly positive. Lean proves this using
\texttt{Real.sqrt\_nonneg} and \texttt{add\_pos\_of\_nonneg\_of\_pos}.
\end{proof}

\begin{theorem}[Advantage sign correctness]
\LeanName{advantage_sign_correct}. If $r_j<\mu<r_i$ and $\epsilon>0$, then
\[
  A_j<0
  \qquad\text{and}\qquad
  0<A_i.
\]
\end{theorem}
\begin{proof}
By \LeanName{groupDenom_pos}, the denominator of each advantage is strictly
positive. The numerator $r_j-\mu$ is negative, so $A_j<0$. The numerator
$r_i-\mu$ is positive, so $A_i>0$. Lean uses
\texttt{div\_neg\_of\_neg\_of\_pos} and \texttt{div\_pos}.
\end{proof}

\begin{theorem}[Winner advantage positivity]
\LeanName{B_winner_advantage_pos}. If an index $i_w$ is a strict group winner
and its reward is above the group mean, and if $\epsilon>0$, then
\[
  0<A_{i_w}.
\]
\end{theorem}
\begin{proof}
The theorem uses the same denominator argument. Since the winner reward is
assumed above the mean, its advantage numerator is positive. Since the
denominator is positive, the quotient is positive. Lean records the strict
winner hypothesis for the intended B-winner interpretation, while the algebraic
sign proof uses the explicit above-mean assumption.
\end{proof}


\section{DPO Directionality Theorems}
\label{sec:dpo-directionality}

All three components of $\RewardTrain$ carry a certified DPO ascent theorem
under the local hypotheses \textbf{H1} (non-degeneracy: the step transfers
strictly positive probability mass $\delta>0$ from loser to winner) and
\textbf{H2} (small step: the transfer is pairwise, other outputs fixed to
first order):
\begin{enumerate}
  \item \LeanName{DPO_B_gradient_correct} (existing; proved above): a DPO
    step on a B-valid pair strictly increases
    $\E_{y\sim\pi_\theta}[\RewardB(q,y)]$.
  \item \LeanName{DPO_C_gradient_correct} (new): a DPO step on a C-valid pair
    strictly increases $\E_{y\sim\pi_\theta}[\RewardC(q,y)]$.
  \item \LeanName{DPO_D_tilde_gradient_correct} (new): a DPO step on a
    $\tilde D$-valid pair strictly increases
    $\E_{y\sim\pi_\theta}[\RewardDT(q,y)]$.
\end{enumerate}

\begin{theorem}[DPO-C gradient correctness]
\LeanName{DPO_C_gradient_correct}. Under H1 and H2, a DPO step on a C-valid
preference pair strictly increases
$\E_{y\sim\pi_\theta}[\RewardC(q,y)]$.
\end{theorem}
\begin{proof}
The proof template is reused from \LeanName{DPO_B_gradient_correct}:
(1) a C-valid pair gives $\RewardC(s_w)>\RewardC(s_l)$ via the utility gate
(\LeanName{C1_utility_gate} dual) and the novelty ordering
(\LeanName{cValidPair_reward_ordered});
(2) \LeanName{dpo_logit_direction} applies unchanged;
(3) \LeanName{expectedReward_pairwiseTransfer_eq} applies unchanged, giving
the expectation change $\delta(\RewardC(s_w)-\RewardC(s_l))>0$;
(4) conclude the strict increase in $\E[\RewardC]$.
\end{proof}

\begin{theorem}[DPO-$\tilde D$ gradient correctness]
\LeanName{DPO_D_tilde_gradient_correct}. Under H1 and H2, a DPO step on a
$\tilde D$-valid preference pair strictly increases
$\E_{y\sim\pi_\theta}[\RewardDT(q,y)]$.
\end{theorem}
\begin{proof}
Identical template to the previous theorem, substituting
\LeanName{D_tilde_dominance} for \LeanName{B1_dominance}. The generic pieces
\LeanName{dpo_logit_direction} and
\LeanName{expectedReward_pairwiseTransfer_eq} are reused unchanged.
\end{proof}


\section{Dense Trajectory Reward and 2-GRPO Quantization}

\subsection{Step-CUE Dense Reward}

The dense trajectory reward $r_t = C(t+1) - C(t)$ is the incremental Step-CUE
contribution of each reasoning step (\LeanName{stepReward}). Under the
per-step utility gate, PROOF-04
(\LeanName{CUE_trajectory_monotone_nondecreasing},
Part~\ref{part:benchmark}) certifies the cumulative signal never punishes a
gated step: the running curve is monotone non-decreasing, so $r_t$ is a
well-posed nonnegative dense reward.

\subsection{DC Bonus Necessity Theorem}

\begin{theorem}[PROOF-15: the terminal DC bonus is necessary]
\LeanName{dense_reward_DC_bonus_incentivizes_exploration}. Consider a
monotone-convergent trajectory and a diverge-converge trajectory with equal
cumulative Step-CUE $c$. Without the bonus their total rewards tie
($\LeanName{totalReward}(c,0) = \LeanName{totalReward}(c,0)$), so the dense
reward alone never prefers exploration; with any strictly positive terminal
bonus $b$ awarded to the diverge-converge trajectory,
\[
  \LeanName{totalReward}(c, 0) < \LeanName{totalReward}(c, b).
\]
Hence without $r_{DC}$ the optimizer is indifferent to (and by tie-breaking
converges on) monotone trajectories with $\mathrm{DC}=0$ despite maximal
cumulative CUE; the bonus strictly breaks the tie toward exploration.
\end{theorem}

\subsection{2-GRPO Advantage Quantization}

\begin{theorem}[PROOF-16: advantages quantize to $\pm 1/\sqrt{2}$]
\LeanName{twoGRPO_advantage_quantization_unbiased}, with
\LeanName{groupScale_eq} and \LeanName{advantage_antisymm}. In 2-GRPO with two
rollouts and rewards $r_1 > r_2$, the group scale collapses to
$\sqrt{(r_1-\mu)^2 + (r_2-\mu)^2} = |r_1 - r_2|/\sqrt{2}$ at
$\mu = (r_1+r_2)/2$, and the group-relative advantages are exactly
\[
  A_1 = \tfrac{1}{\sqrt{2}},
  \qquad
  A_2 = -\tfrac{1}{\sqrt{2}},
\]
independent of the reward magnitudes. The advantage is therefore a
deterministic function of the signed rank alone --- an unbiased signed-rank
estimator, $\E[A_i \mid \operatorname{rank}(y_i)] = \pm 1/\sqrt{2}$ --- which
is the formal core of the GRPO-is-secretly-DPO reduction used by the two-phase
2-CGRPO protocol.
\end{theorem}

\section{Goodhart-Resistance Conditions}
\label{sec:goodhart-conditions}

A 2-CGRPO run is valid only if all four of the following axioms hold as
explicit training preconditions. Each is stated in Lean and collected in the
summary conjunction below.

\begin{enumerate}
  \item \LeanName{D_tilde_no_fixed_target} --- $Q_{\mathrm{probe}}$ is
    resampled every batch; no probe string is fixed across training steps.
  \item \LeanName{C_corpus_stationarity_condition} --- the corpus
    $\mathcal{H}$ used to compute $N_{\mathrm{KL}}$ is updated on a fixed
    external schedule only, never by inserting policy outputs.
  \item \LeanName{B_utility_gate_tightness} --- $\tau$ is calibrated above
    the entropy of incoherent baseline outputs.
  \item \LeanName{ProxyFaithfulness_decay_detectable} ---
    $\rho(\RewardDT,\RewardD)$ is monitored at every checkpoint; training
    halts if $\rho<\rho_{\min}$ (Part~\ref{part:proxy}).
\end{enumerate}

\begin{assumption}[Corpus stationarity for C]
\LeanName{C_corpus_stationarity_condition}. The corpus $\mathcal{H}$ used to
compute $N_{\mathrm{KL}}$ is updated on a fixed external schedule, never by
inserting policy outputs: the corpus trajectory equals a schedule fixed
before training, and no policy output produced during training occurs inside
the corpus it is scored against (\LeanName{CorpusStationarity}). Violating
this axiom breaks the anti-copy and anti-padding guarantees
(\LeanName{C2_copy_penalty}, \LeanName{C3_no_padding_exploit}) by making
$\mathcal{H}$ self-referential: the policy would control the compression
baseline that defines its own novelty. This is the formal Goodhart-resistance
condition for $\RewardC$.
\end{assumption}

\begin{assumption}[Utility gate tightness for B]
\LeanName{B_utility_gate_tightness}. The existing B theorems prove
anti-collapse but do not block entropy-maximizing \emph{incoherence}:
outputs achieving high $H_q(y)$ by being vague rather than expansive. The
utility threshold must satisfy
\[
  \tau > H_q(y_{\mathrm{random}})
\]
for every feasibility-passing but semantically incoherent output
$y_{\mathrm{random}}$ (\LeanName{IncoherentBaseline}). This is a calibration
requirement: $\tau$ must be set above the entropy of incoherent baseline
outputs, not just above zero, and the calibration must be validated against
incoherent high-entropy baselines \emph{before training}. Without it,
$\RewardB$ can be Goodharted by entropy-maximizing incoherence. The Lean
consequences \LeanName{B_incoherence_reward_below_threshold} and
\LeanName{B_coherent_dominates_incoherent} show that under a tight gate no
incoherent baseline can reach the calibrated bar or dominate a genuinely
expansive output.
\end{assumption}

\begin{theorem}[Goodhart resistance summary]
\LeanName{Goodhart_resistance_summary}. For any admissible protocol
certificates, the four Goodhart conditions hold simultaneously:
\[
  \text{\LeanName{D_tilde_no_fixed_target}}
  \;\wedge\;
  \text{\LeanName{C_corpus_stationarity_condition}}
  \;\wedge\;
  \text{\LeanName{B_utility_gate_tightness}}
  \;\wedge\;
  \text{\LeanName{ProxyFaithfulness_decay_detectable}}.
\]
\end{theorem}
\begin{proof}
Lean assembles the conjunction from the four protocol certificates
(\LeanName{DTildeTrainingProtocol}, \LeanName{CTrainingProtocol},
\LeanName{TauCalibration}, \LeanName{ProxyMonitor}), so the full protection
is checkable in one place.
\end{proof}

\noindent
\textbf{Validity statement.} A training run that violates any of these four
conditions does not have Lean-certified Goodhart resistance, and its results
cannot be compared against the benchmark metrics of
Part~\ref{part:benchmark}.


%======================================================================
\part{Proxy Validation and Resolution Conditions}
\label{part:proxy}
%======================================================================

\section{Proxy Faithfulness for D}
\label{sec:proxy-faithfulness}

Training ascends $\RewardDT$ on resampled probe prompts; the benchmark
measures $\RewardD$ on the frozen probe set $\mathcal{P}$. The bridge between
them is conditional on an empirical assumption.

\begin{assumption}[Proxy faithfulness]
\LeanName{ProxyFaithfulness}. After monotone calibration, the KL shift on
resampled $q_{\mathrm{probe}}$ and $\RewardD$ on the frozen $\mathcal{P}$ are
monotonically related up to error bound $\delta$: uniformly over outputs,
\[
  \bigl|\RewardDT(s)-\RewardD(s,\mathcal{H},\mathcal{P})\bigr|\le\epsilon.
\]
Unlike the purely algebraic \LeanName{BEstimator}/\LeanName{CEstimator}
bounds, this is an \emph{empirical} assumption that the benchmark correlation
study must estimate before training begins.
\end{assumption}

\begin{theorem}[Proxy rank agreement]
\LeanName{D_tilde_proxy_rank_agreement}. Under \LeanName{ProxyFaithfulness},
if
\[
  \RewardDT(s_1)-\RewardDT(s_2)>2\epsilon,
\]
then
\[
  \RewardD(s_1,\mathcal{H},\mathcal{P})>\RewardD(s_2,\mathcal{H},\mathcal{P}).
\]
\end{theorem}
\begin{proof}
Structure mirrors \LeanName{B_rank_preserved} and
\LeanName{C_rank_preserved}: extract the one-sided bounds from the uniform
error bound with \texttt{abs\_le.mp}, then chain the inequalities with
\texttt{linarith}.
\end{proof}

\begin{theorem}[Proxy faithfulness decay is detectable]
\LeanName{ProxyFaithfulness_decay_detectable}. Let a
\LeanName{ProxyMonitor} record the empirical correlation
$\rho(\RewardDT,\RewardD)$ at every training checkpoint, with threshold
$\rho_{\min}$ and the calibration link that faithfulness implies
$\rho\ge\rho_{\min}$. If at checkpoint $t$ the measured correlation falls
below the threshold, $\rho(t)<\rho_{\min}$, then proxy faithfulness fails at
$t$: the rank-agreement guarantee of
\LeanName{D_tilde_proxy_rank_agreement} no longer holds, and training must
halt or the proxy must be recalibrated.
\end{theorem}
\begin{proof}
If faithfulness held at $t$, the calibration link would force
$\rho_{\min}\le\rho(t)$, contradicting the decay hypothesis. Lean closes with
\texttt{not\_le}.
\end{proof}

\noindent
\textbf{Calibration requirement.} $\rho_{\min}$ is a parameter of the
monitor: the benchmark correlation study must estimate it \emph{before
training begins}. This theorem is the \textbf{formal Goodhart early-warning
condition}: it converts the proxy faithfulness assumption into a monitorable
training health invariant rather than a one-time post-hoc validation.

\section{Training-to-Benchmark Correspondence}
\label{sec:correspondence-table}

\begin{table}[h]
\centering
\small
\begin{tabular}{@{}p{0.20\linewidth}p{0.20\linewidth}p{0.28\linewidth}p{0.20\linewidth}@{}}
\toprule
Training metric & Benchmark metric & Correspondence theorem & Resolution condition \\
\midrule
$\RewardB$ &
$\widehat{H}_q(y)$ on held-out tasks &
\LeanName{B_rank_preserved} &
Gap $>2\epsilon$ \\
$\RewardC$ via $N_{\mathrm{KL}}$ (Composability Score only) &
$\widehat{N}(s\mid\mathcal{H})$ on held-out corpus &
\LeanName{C_rank_preserved} &
Gap $>2\epsilon$; not part of Benchmark Score \\
$\RewardDT$ on resampled $q_{\mathrm{probe}}$ &
$\RewardD$ on frozen $\mathcal{P}$ &
\LeanName{D_tilde_proxy_rank_agreement} &
Gap $>2\epsilon$ and $\rho\ge\rho_{\min}$ \\
None (by design) &
$G_k$ interaction gain &
Detection when gap $>2\epsilon$
(\LeanName{MAS_benchmark_detects_irreducibility},
\LeanName{kAgent_benchmark_detects}) &
Pure benchmark --- no training correspondent \\
\bottomrule
\end{tabular}
\caption{One-to-one pairing of training metrics with benchmark metrics, the
Lean correspondence theorem, and the resolution condition under which the
correspondence is valid.}
\end{table}

\begin{remark}[The one asymmetry]
The $\rho_{\min}$ condition in the $\RewardDT$ row is the \textbf{only
asymmetry} among the training--benchmark pairs. B has purely algebraic rank
preservation; D's correspondence additionally requires empirical proxy
faithfulness. This asymmetry is the formal price of making $\RewardD$
benchmark-only. The reported Benchmark Score is~\eqref{eq:Rcreativity}, not
an additive $B$+$D$ sum; training uses
$\RewardTrain=\RewardB+\tilde\lambda_D\RewardDT$.
\end{remark}


%======================================================================
\part{The 2-CGRPO Algorithm}
\label{part:algorithm}
%======================================================================

\section{Two-Phase Protocol}
\label{sec:two-phase}

\subsection*{Training phase}
\begin{itemize}
  \item Compute $\RewardTrain=\RewardB+\tilde\lambda_D\RewardDT$ per output,
    with a fresh $q_{\mathrm{probe}}$ sample per batch (no $R_C$ channel).
  \item Normalize to group-relative advantages via
    \LeanName{groupDenom_pos}, \LeanName{advantage_sign_correct}, and
    \LeanName{trainReward_advantage_sign_correct}.
  \item Update the policy via DPO/GRPO; the directionality certificates are
    the B and $\tilde D$ theorems of Section~\ref{sec:dpo-directionality}.
  \item At every checkpoint: compute $\rho(\RewardDT,\RewardD)$ on a small
    held-aside validation set; if $\rho<\rho_{\min}$, halt and flag a
    \LeanName{ProxyFaithfulness} violation
    (\LeanName{ProxyFaithfulness_decay_detectable}).
\end{itemize}

\subsection*{Evaluation phase}
\begin{itemize}
  \item Freeze the policy.
  \item Compute the Benchmark Score~\eqref{eq:Rcreativity}
    (\LeanName{Rcreativity}) on held-out tasks with the frozen probe set
    $\mathcal{P}$ (Section~\ref{sec:benchmark-score}).
  \item Compute $G_k$ across agent configurations from ContribBench.
  \item Check all resolution conditions of
    Section~\ref{sec:correspondence-table} (gap $>2\epsilon$ for each
    metric).
  \item Report the \LeanName{ProxyFaithfulness} empirical correlation as a
    first-class result --- not an appendix detail.
\end{itemize}

\paragraph{Naming.}
\textbf{2-CGRPO} denotes: two-phase (train, then certify), Creative
objective, Group Relative Policy Optimization. The ``2'' also signals the
two-layer protection: algebraic Lean guarantees at training time, empirical
proxy validation at benchmark time.

\section{Validation Ledger}

\begin{table}[h]
\centering
\scriptsize
\begin{tabular}{@{}p{0.18\linewidth}p{0.34\linewidth}p{0.40\linewidth}@{}}
\toprule
Layer & Lean theorem & Main assumptions \\
\midrule
Entropy & \LeanName{shannonEntropy_nonneg} & finite support PMF \\
Entropy & \LeanName{shannonEntropy_point_mass} & finite support PMF, point mass \\
Entropy & \LeanName{shannonEntropy_lt_iff} & definitional unfolding \\
B helper & \LeanName{rewardB_of_feasible} & feasible output \\
B helper & \LeanName{rewardB_of_infeasible} & infeasible output \\
B helper & \LeanName{rewardB_nonneg} & entropy nonnegativity and gate cases \\
B helper & \LeanName{mem_feasibleSet_iff} & finite filtered feasible set \\
B objective & \LeanName{B1_dominance} & feasible pair, entropy gap \\
B objective & \LeanName{B2_no_collapse} & finite $Y,Z$, feasible nonempty set, entropy gap \\
B mixture & \LeanName{B3_mixture_lb} & mass condition $\RewardB(y_{\mathrm{low}})\le\pi(y_{\mathrm{high}})\RewardB(y_{\mathrm{high}})$ \\
B mixture & \LeanName{entropy_mixture_concavity} & \textbf{theorem} (PROOF-09, promoted): concave Jensen on \texttt{negMulLog} \\
B mixture & \LeanName{B3b_jensen_anti_collapse} & now axiom-free via \LeanName{entropy_mixture_concavity} \\
C objective & \LeanName{C1_utility_gate} & $u<\tau$ (Composability Score only) \\
C objective & \LeanName{C2_copy_penalty} & copy-minimal compressor (Composability Score only) \\
C objective & \LeanName{C3_no_padding_exploit} & normalized compressor padding bound (Composability Score only) \\
C objective & \LeanName{C3_no_trivial_padding_gain} & zero-padding cost, normalized compressor (Composability Score only) \\
N surface & \LeanName{N1_utility_gate} & $u<\tau$ \\
N surface & \LeanName{N2_kl_nonneg} & Gibbs inequality (finite PMF axiom) \\
N surface & \LeanName{N3_kl_component_bound} & KL plus entropy bounded by novelty \\
N surface & \LeanName{klDivergence_nonneg} & finite PMF KL nonnegativity axiom \\
N surface & \LeanName{kl_cross_entropy_decomposition} & cross-entropy link to C \\
D probe-relative & \LeanName{D1_utility_gate} & $u<\tau$, fixed finite $\mathcal{P}$ \\
D probe-relative & \LeanName{D2_dominance} & feasibility, strict deformation gap \\
D probe-relative & \LeanName{D3_rewardC_bounded} & copy-minimal compressor (Composability Score only) \\
D probe-relative & \LeanName{D3_exploration_transformation_bound} & \LeanName{ProbeCompressor} (Composability Score only) \\
D probe-relative & \LeanName{D3_rewardD_bound_standalone} & \LeanName{ProbeCompressor}; Benchmark Score, $R_C$-free \\
D probe-relative & \LeanName{D4_zero_for_corpus_copy} & append-self invariance \\
D probe-relative & \LeanName{D4_zero_reward_for_corpus_copy} & feasibility plus D4 \\
B baseline & \LeanName{BBase_indistinguishable} & feasible pair, equal utility \\
B baseline & \LeanName{BContrast_distinguishes} & feasible pair, equal utility, entropy gap \\
C baseline & \LeanName{CBase_infeasible_can_maximize} & $u<\tau$, positive novelty (Composability Score only) \\
B estimator & \LeanName{B_rank_preserved} & uniform error, true gap $>2\epsilon$ \\
C estimator & \LeanName{C_rank_preserved} & uniform error, true gap $>2\epsilon$ \\
DPO bridge & \LeanName{dpo_logit_direction} & distinct winner/loser, $\eta>0$ \\
DPO bridge & \LeanName{DPO_B_correct_direction} & feasible B pair, entropy gap, $\eta>0$ \\
DPO bridge & \LeanName{dpo_gradient_increases_logodds} & $\beta>0$, positive sigmoid term, H1, $\eta>0$ \\
DPO bridge & \LeanName{softmax_monotone_in_logit_diff} & two-logit softmax, $\delta>0$ \\
DPO bridge & \LeanName{dpo_step_increases_expected_RB} & B-valid pair, positive probability transfer \\
DPO bridge & \LeanName{DPO_B_gradient_correct} & local DPO transfer hypothesis (H1/H2) \\
RL signal & \LeanName{benchmarkReward_nonneg} & nonnegative B/D rewards and $\lambda_D$ (diagnostic) \\
RL signal & \LeanName{trainReward_nonneg} & nonnegative B/$\tilde D$ training signal (no $R_C$) \\
RL signal & \LeanName{trainReward_neq_benchmarkReward} & train/benchmark split machine-checked \\
RL signal & \LeanName{groupDenom_pos} & positive advantage stabilizer $\epsilon$ \\
RL signal & \LeanName{advantage_sign_correct} & rewards below/above group mean \\
RL signal & \LeanName{B_winner_advantage_pos} & winner above group mean, $\epsilon>0$ \\
\bottomrule
\end{tabular}
\caption{Lean-checked claims and their principal assumptions (objective,
estimator, and DPO layers).}
\end{table}

\begin{table}[ht]
\centering
\scriptsize
\begin{tabular}{@{}p{0.18\linewidth}p{0.34\linewidth}p{0.40\linewidth}@{}}
\toprule
Layer & Lean theorem & Main assumptions \\
\midrule
MAS & \LeanName{MAS_irreducibility} & positive interaction gain \\
MAS benchmark & \LeanName{MAS_benchmark_detects_irreducibility} & interaction gain $>2\epsilon$ \\
MAS $k$-agent & \LeanName{kAgent_irreducibility} & positive $k$-agent interaction gain \\
MAS $k$-agent & \LeanName{kAgent_benchmark_detects} & $k$-agent gain $>2\epsilon$ \\
MAS graph & \LeanName{CommGraph} & irreflexive directed communication graph \\
MAS graph & \LeanName{nonProduct_base_case} & faithfulness plus edge $0\to1$ \\
MAS graph & \LeanName{connectedMAS_nonProductKernel} & faithfulness plus at least one edge \\
MAS graph & \LeanName{gain_zero_iff_joint_eq_max} & \textbf{theorem} (PROOF-11, corrected): zero gain iff $H_J=\max_i H_i$; old axiom removed \\
MAS graph & \LeanName{productKernel_additivity_gain_pos} & additive joint entropy, positive marginals $\Rightarrow$ positive gain \\
MAS graph & \LeanName{gain_biconditional} & nonnegative gain plus corrected degeneracy characterization \\
O-info & \LeanName{oInformation3_negative_of_pair_sum_lt} & pair-sum below triple term \\
O-info & \LeanName{oInformation3_positive_of_triple_lt_pair_sum} & triple term below pair-sum \\
O-info & \LeanName{equalMarginals_gain_implies_synergy} & equal-marginal sign condition \\
O-info & \LeanName{negativeGain_implies_redundancy} & collapse sign condition \\
O-info bridge & \LeanName{productKernel_joint_entropy_eq_sum} & \textbf{theorem} (PROOF-10, promoted): via \LeanName{shannonEntropy_jointPMF} \\
O-info bridge & \LeanName{productKernel_implies_redundancy} & strictly positive marginal entropies \\
Composability & \LeanName{Bmax_not_jointly_dominated} & B maximizer \\
Composability & \LeanName{paretoFrontier_nonempty} & finite nonempty output type \\
Composability & \LeanName{BC_tension_example} & explicit $\operatorname{Fin}(2)$ construction \\
Composability & \LeanName{BCD_tension_example} & explicit $\operatorname{Fin}(3)$ construction \\
Composability & \LeanName{BCD_paretoFrontier_nonempty} & finite nonempty output type \\
Composability & \LeanName{paretoFrontier_nonempty_BD} & finite nonempty output type; Benchmark Score, $R_C$-free \\
Estimator calibration & \LeanName{plugIn_bias_axiom} & \textbf{registered import} (PROOF-12, KG-2): Miller--Madow-scale bound \\
Estimator calibration & \LeanName{entropy_ranking_sample_requirement} & true gap $>|Y|/n$ \\
Benchmark spec & \LeanName{benchmark_can_detect_S2_dominance} & true S2--S1 gap $>4\epsilon$ \\
\bottomrule
\end{tabular}
\caption{Lean-checked claims and their principal assumptions (multi-agent,
composability, and benchmark layers).}
\end{table}

\begin{table}[ht]
\centering
\scriptsize
\begin{tabular}{@{}p{0.18\linewidth}p{0.34\linewidth}p{0.40\linewidth}@{}}
\toprule
Layer & Lean theorem & Main assumptions \\
\midrule
D proxy & \LeanName{D_tilde_utility_gate} & $u<\tau$ \\
D proxy & \LeanName{D_tilde_nonneg} & gated KL nonnegativity \\
D proxy & \LeanName{D_tilde_zero_for_identical} & policy matches frozen reference on probe \\
D proxy & \LeanName{D_tilde_no_fixed_target} & \textbf{protocol axiom}: probe batch resampled, no fixed probe string \\
D proxy & \LeanName{D_tilde_dominance} & feasibility, strict proxy-reward gap \\
DPO bridge & \LeanName{DPO_C_gradient_correct} & C-valid pair, local DPO transfer (H1/H2) \\
DPO bridge & \LeanName{DPO_D_tilde_gradient_correct} & $\tilde D$-valid pair, local DPO transfer (H1/H2) \\
Proxy bridge & \LeanName{D_tilde_proxy_rank_agreement} & proxy faithfulness, proxy gap $>2\epsilon$ \\
Proxy bridge & \LeanName{ProxyFaithfulness_decay_detectable} & monitor calibration link, $\rho<\rho_{\min}$ \\
RL signal & \LeanName{trainReward_nonneg} & nonnegative B/C components and weights \\
RL signal & \LeanName{trainReward_advantage_sign_correct} & train rewards below/above group mean \\
RL signal & \LeanName{trainReward_neq_fullReward} & weighted D components differ \\
Goodhart & \LeanName{B_utility_gate_tightness} & \textbf{calibration axiom}: $\tau>H_q(y_{\mathrm{random}})$ \\
Goodhart & \LeanName{C_corpus_stationarity_condition} & \textbf{protocol axiom}: external corpus schedule only \\
Goodhart & \LeanName{Goodhart_resistance_summary} & all four protocol certificates \\
\bottomrule
\end{tabular}
\caption{Lean-checked claims and their principal assumptions (training-signal,
proxy, and Goodhart layers).}
\end{table}

\subsection*{CUE-Layer and Structural Additions}

\begin{table}[h]
\centering
\scriptsize
\begin{tabular}{@{}p{0.14\linewidth}p{0.36\linewidth}p{0.28\linewidth}p{0.14\linewidth}@{}}
\toprule
Layer & Lean theorem & Main assumptions & Status \\
\midrule
Foundation & \LeanName{maxEntropy_finite_alphabet} & Gibbs inequality (\LeanName{klDivergence_nonneg}) & Machine-checked \\
Foundation & \LeanName{klDivergence_uniform_eq} & finite nonempty alphabet & Machine-checked \\
Foundation & \LeanName{sum_toReal_eq_one} & finite PMF & Machine-checked \\
Foundation & \LeanName{shannonEntropy_eq_sum_negMulLog} & definitional & Machine-checked \\
CUE & \LeanName{cue}, \LeanName{cue_nonneg} & nonnegative Brier delta, positive bit length & Machine-checked \\
CUE & \LeanName{CUE_voi_equivalence} (PROOF-01) & proper-scoring-rule bound field & Machine-checked \\
CUE & \LeanName{CUE_voi_equality} (PROOF-01) & calibrated binary saturation & Machine-checked \\
CUE & \LeanName{shannon_upper_bound_normalization} (PROOF-02) & VOI $\le$ posterior entropy field & Machine-checked \\
CUE & \LeanName{CUE_primary_gate} & gate definition (Composability Score only; superseded by shifted form) & Machine-checked \\
CUE & \LeanName{receiver_expansion_clarification_separation} (PROOF-06) & explicit $\operatorname{Fin}(2)$ construction & Machine-checked \\
CUE shifted & \LeanName{CUE_shifted_bound} & Shannon ceiling minus $f$; Benchmark Score, $R_C$-free & Machine-checked \\
CUE shifted & \LeanName{CUE_shifted_sign_iff} & threshold clearance; Benchmark Score, $R_C$-free & Machine-checked \\
CUE shifted & \LeanName{receiver_expansion_bounded} & normalized entropy ratio; Benchmark Score, $R_C$-free & Machine-checked \\
CUE shifted & \LeanName{Rcreativity_sign_matches_CUE_shifted} & $\alpha>0$, $r_b\in[0,1]$; Benchmark Score, $R_C$-free & Machine-checked \\
CUE shifted & \LeanName{Gk_gated_by_CUE_shifted} & bridge implication, $f<\mathrm{CUE}$; Benchmark Score, $R_C$-free & Machine-checked \\
Trajectory & \LeanName{CUE_trajectory_monotone_nondecreasing} (PROOF-04) & per-step gate $r_t\ge0$ & Machine-checked \\
Trajectory & \LeanName{DC_score_existence} (PROOF-05) & explicit tent construction & Machine-checked \\
Trajectory & \LeanName{monotone_increasing_not_DC} (PROOF-05) & strictly increasing trajectory & Machine-checked \\
Trajectory & \LeanName{step_CUE_gamma_positivity_characterization} (PROOF-14) & fitted exponential-plus-linear curve & Machine-checked \\
Trajectory & \LeanName{dense_reward_DC_bonus_incentivizes_exploration} (PROOF-15) & equal cumulative CUE, positive bonus & Machine-checked \\
\bottomrule
\end{tabular}
\caption{CUE-layer additions (foundation, gate, and trajectory diagnostics)
with verification status. ``Machine-checked'' means proved in Lean with no
\texttt{sorry}.}
\end{table}

\begin{table}[ht]
\centering
\scriptsize
\begin{tabular}{@{}p{0.14\linewidth}p{0.36\linewidth}p{0.28\linewidth}p{0.14\linewidth}@{}}
\toprule
Layer & Lean theorem & Main assumptions & Status \\
\midrule
MAS bridge & \LeanName{Gk_implies_CUE_improvement} (PROOF-07) & calibration link, gain $>2\varepsilon_{\mathrm{cal}}$ & Machine-checked \\
MAS bridge & \LeanName{Gk_implies_CUE_improvement_sup} (PROOF-07) & calibration link, gain $>2\varepsilon_{\mathrm{cal}}$ & Machine-checked \\
DPS & \LeanName{dual_process_separation_principle} (PROOF-03) & live-only reward construction & Machine-checked \\
DPS & \LeanName{Rtrain_cannot_goodhart_frozen} (PROOF-03) & live-only reward construction & Machine-checked \\
Goodhart & \LeanName{rank_reversal_refutes_agreement} (PROOF-13) & rank-reversal witness & Machine-checked \\
Goodhart & \LeanName{reversal_instance_exists} (PROOF-13) & explicit two-candidate construction & Machine-checked \\
Goodhart & \LeanName{structural_goodhart_collapse} (PROOF-13) & composition of the above & Machine-checked \\
Commensurability & \LeanName{rewardB_le_maxEntropy} (PROOF-08) & max-entropy bound & Machine-checked \\
Commensurability & \LeanName{lambda_normalization_equalizes_maxima} (PROOF-08) & positive length and $\log|\Sigma|$ (Composability Score only) & Machine-checked \\
Commensurability & \LeanName{lambda_D_normalization_standalone} & probe size/length only; Benchmark Score, $R_C$-free & Machine-checked \\
2-GRPO & \LeanName{groupScale_eq} (PROOF-16) & two-rollout group & Machine-checked \\
2-GRPO & \LeanName{advantage_antisymm} (PROOF-16) & two-rollout group & Machine-checked \\
2-GRPO & \LeanName{twoGRPO_advantage_quantization_unbiased} (PROOF-16) & $r_2<r_1$ & Machine-checked \\
MAS promoted & \LeanName{jointPMF_apply} (PROOF-10) & finite alphabet & Machine-checked \\
MAS promoted & \LeanName{shannonEntropy_jointPMF} (PROOF-10) & finite alphabet & Machine-checked \\
Known gap & \LeanName{klDivergence_nonneg} (KG-1) & Gibbs inequality & Registered import \\
Known gap & \LeanName{kl_cross_entropy_decomposition} (KG-1) & cross-entropy link & Registered import \\
Known gap & \LeanName{plugIn_bias_axiom} (KG-2) & Miller--Madow bound & Registered import \\
\bottomrule
\end{tabular}
\caption{Structural additions (dual-process separation, Goodhart collapse,
commensurability, 2-GRPO, MAS promotions) with verification status.
``Registered import'' means an explicitly documented trusted axiom listed in
the Known Gaps Registry (Appendix~\ref{app:known-gaps}).}
\end{table}

\appendix

\section{Benchmark Hold-Out Probe Protocol}
\label{app:probe-protocol}

Component D is undefined until the probe set $\mathcal{P}$ is fixed. The
Component D section defines the hold-out protocol; experiments report the
concrete instantiation here: sampling seed, split index, $|\mathcal{P}|$, and
per-stratum counts. Under this protocol, $\mathcal{P}$ is reproducible and
domain-representative by construction. Lean proves properties of
$D(s,\mathcal{H},\mathcal{P})$ for any fixed finite $\mathcal{P}$; the benchmark
protocol fixes which finite $\mathcal{P}$ is used.

\section{Computational Approximation: Surface KL Novelty}
\label{app:n-surface}

The file \texttt{Creativity/C/SurfaceNovelty.lean} defines surface
distributional novelty on finite PMFs. For output distribution $p_s$ and corpus
distribution $p_{\mathcal{H}}$,
\[
  N_{\mathrm{KL}}(s\mid\mathcal{H})
  :=
  D_{\mathrm{KL}}(p_s\|p_{\mathcal{H}}),
\]
formalized as \LeanName{surfaceNoveltyKL}. This is an empirical approximation
to C, not an independent creativity dimension.

\begin{theorem}[N1: utility gate]
\LeanName{N1_utility_gate}. If $u<\tau$, then $\RewardN(q,s)=0$.
\end{theorem}
\begin{proof}
Identical in structure to C1: infeasibility forces the zero branch of the gated
reward definition. Lean simplifies \texttt{rewardN} with \texttt{not\_le.mpr}.
\end{proof}

\begin{theorem}[N2: KL nonnegativity]
\LeanName{N2_kl_nonneg}. For all finite-alphabet PMFs $p_s,p_{\mathcal{H}}$,
\[
  0\le D_{\mathrm{KL}}(p_s\|p_{\mathcal{H}}).
\]
\end{theorem}
\begin{proof}
This is Gibbs' inequality for finite distributions. Lean invokes the axiom
\LeanName{klDivergence_nonneg} from
\texttt{Creativity/Probability/KLDivergence.lean}.
\end{proof}

\begin{theorem}[N3: KL as a component of algorithmic novelty]
\LeanName{N3_kl_component_bound}. If
\[
  D_{\mathrm{KL}}(p_s\|p_{\mathcal{H}})+\Ent(p_s)\le N(s\mid\mathcal{H}),
\]
then
\[
  D_{\mathrm{KL}}(p_s\|p_{\mathcal{H}})\le N(s\mid\mathcal{H}).
\]
\end{theorem}
\begin{proof}
Both $D_{\mathrm{KL}}$ and $\Ent(p_s)$ are nonnegative by N2 and
\LeanName{shannonEntropy_nonneg}. Subtracting the entropy term from the assumed
bound yields the claim. Lean closes with \texttt{linarith}. The cross-entropy
decomposition \LeanName{kl_cross_entropy_decomposition} axiomatizes the precise
link $N_{\mathrm{KL}}+\Ent(p_s)\approx N(s\mid\mathcal{H})$ when strings are
mapped to PMFs.
\end{proof}

\section{Composability Scaffolding}
\label{app:composability-scaffolding}

The main text uses the Pareto-frontier theorem. The following finite
constructions are Lean sanity checks showing that the objectives can prefer
different outputs.

\begin{theorem}[B-maximizer is not jointly dominated]
\LeanName{Bmax_not_jointly_dominated}. If $y_{\max}$ satisfies
\[
  B(y)\le B(y_{\max})\qquad\text{for all }y,
\]
then there is no $y'$ such that
\[
  B(y_{\max})<B(y')
  \quad\text{and}\quad
  C(y_{\max})\le C(y').
\]
\end{theorem}
\begin{proof}
Any joint dominator would have to be strictly better in B:
$B(y_{\max})<B(y')$. But the maximality assumption gives
$B(y')\le B(y_{\max})$, contradicting strict inequality. Lean uses
\texttt{not\_lt.mpr}.
\end{proof}

\begin{theorem}[B--C tension example]
\LeanName{BC_tension_example}. There exist rewards $B,C:\operatorname{Fin}(2)\to
\R$ such that output $0$ maximizes B, output $1$ maximizes C, and the preferences
are genuinely opposed:
\[
  B(1)<B(0),\qquad C(0)<C(1).
\]
\end{theorem}
\begin{proof}
Take $B(0)=1$, $B(1)=0$, $C(0)=0$, and $C(1)=1$. Then $0$ is the B-maximizer,
$1$ is the C-maximizer, and the two strict inequalities hold immediately. Lean
proves the finite cases by case analysis on \(\operatorname{Fin}(2)\) and
numeric simplification.
\end{proof}

\begin{theorem}[B--C--D tension example]
\LeanName{BCD_tension_example}. There exist rewards
$B,C,D:\operatorname{Fin}(3)\to\R$ such that output $0$ maximizes B, output $1$
maximizes C, output $2$ maximizes D, and the three preferences are genuinely
opposed across dimensions.
\end{theorem}
\begin{proof}
Lean constructs explicit rewards on $\operatorname{Fin}(3)$ and verifies by
finite case analysis that each component is maximized at a different index while
the opposing strict inequalities hold.
\end{proof}

\begin{theorem}[Nonempty three-dimensional Pareto frontier]
\LeanName{BCD_paretoFrontier_nonempty}. For any finite nonempty output type $Y$
and rewards $B,C,D:Y\to\R$, there exists $y:Y$ not jointly dominated in all
three components.
\end{theorem}
\begin{proof}
Same finite-maximizer argument as the two-component case, applied to the B
coordinate while ignoring the C and D inequalities in the domination witness.
Lean uses \texttt{Finset.exists\_mem\_eq\_sup'} and \texttt{not\_lt.mpr}.
\end{proof}

\section{Lean Appendix: B3 Scaffolding}
\label{app:b3-scaffolding}

The main text presents \LeanName{B3_mixture_lb} as the general mixture
anti-collapse result. Lean proves it by reduction to the algebraic lemmas below.
These are scaffolding, not independent contributions.

\begin{theorem}[Entropy concavity on downstream mixtures (promoted, PROOF-09)]
\LeanName{entropy_mixture_concavity}. For finite $Y,Z$, model $m$, prompt
$q$, and mixture $\pi:\PMF(Y)$, let $\mu_q(\pi)$ denote the downstream mixture
kernel obtained by sampling $y\sim\pi$ and then $z\sim T(\cdot\mid q,y)$. Then
\[
  \sum_y \pi(y)H_q(y)
  \le
  \Ent(\mu_q(\pi)).
\]
This is Shannon concavity for finite mixtures. Formerly the axiom
\LeanName{entropy_mixture_concavity_axiom}; now a derived theorem via concave
Jensen on \texttt{Real.negMulLog} (see the proof in the B mixture section).
\end{theorem}

\begin{theorem}[B3 scaffolding: weighted mixture lower bound]
\LeanName{B3_weighted_high_lb}. Assume $Y$ and $Z$ are finite. Let
$\pi : \PMF(Y)$ and $y_{\mathrm{low}},y_{\mathrm{high}}:Y$. If
\[
  \RewardB(q,y_{\mathrm{low}})
    \le \pi(y_{\mathrm{high}})\RewardB(q,y_{\mathrm{high}}),
\]
then
\[
  \RewardB(q,y_{\mathrm{low}})
    \le \sum_{y:Y} \pi(y)\RewardB(q,y).
\]
\end{theorem}
\begin{proof}
Every summand in the expectation is nonnegative, so the full finite sum is at
least $\pi(y_{\mathrm{high}})\RewardB(q,y_{\mathrm{high}})$. Composing with the
hypothesis gives the claim. Lean uses \texttt{Finset.sum\_nonneg}.
\end{proof}

\begin{theorem}[B3 scaffolding: strict expected dominance]
\LeanName{B3_expected_exceeds_dominated}. Assume $Y$ and $Z$ are finite. If
$y_{\mathrm{low}}$ and $y_{\mathrm{high}}$ are feasible,
$H_q(y_{\mathrm{low}}) < H_q(y_{\mathrm{high}})$, $\pi(y_{\mathrm{high}})>0$,
and
\[
  \RewardB(q,y_{\mathrm{low}})
    < \pi(y_{\mathrm{high}})\RewardB(q,y_{\mathrm{high}}),
\]
then
\[
  \RewardB(q,y_{\mathrm{low}})
    < \sum_{y:Y}\pi(y)\RewardB(q,y).
\]
\end{theorem}
\begin{proof}
As above, $\pi(y_{\mathrm{high}})\RewardB(q,y_{\mathrm{high}})\le\sum_y\pi(y)\RewardB(q,y)$.
Transitivity with the strict hypothesis gives the result. Lean uses
\texttt{lt\_of\_lt\_of\_le}.
\end{proof}

\section{Experimental Falsifiability Map}
\label{app:falsifiability}

Each experimental prediction E1--E5 of the proposal corresponds to a formal
statement derivable from the Lean theorems. A failed experiment therefore
falsifies an identified assumption, not the framework diffusely.

\begin{center}
\begin{tabular}{@{}p{0.30\linewidth}p{0.62\linewidth}@{}}
\toprule
Experiment & Formal correspondent \\
\midrule
E1 (Judge calibration) & \LeanName{B_rank_preserved}, \LeanName{C_rank_preserved}: rankings are estimator-noise-proof exactly when the measured gap exceeds $2\epsilon$; judge miscalibration manifests as $\epsilon$ inflation. \\
E2 (Calibration failure replication) & \LeanName{BContrast_distinguishes}: CUE-gated B maintains rank where utility-only reward models are provably indistinguishable. \\
E3 (Repetition vs.\ score) & \LeanName{receiver_expansion_clarification_separation} (PROOF-06): the $R_B^{\to A}$ label is anti-correlated with surface similarity by construction, since equal-CUE outputs carry opposite expansion labels. \\
E4 (Discriminant validity) & \LeanName{BCD_tension_example}, \LeanName{BCD_paretoFrontier_nonempty}: the Pareto frontier is non-empty with distinct maximizers, so the dimensions are provably non-redundant. \\
E5 (Receiver stability) & \LeanName{DC_score_existence} (PROOF-05) $+$ \LeanName{D_tilde_proxy_rank_agreement}: trajectory shape and proxy rank agreement are stable under the resolution conditions. \\
\bottomrule
\end{tabular}
\end{center}

\section{2-GRPO Training Protocol Specification}
\label{app:2grpo-protocol}

This appendix makes the Lean-certified guarantees operationally executable for
the planned training run. The formal object instantiated is
\LeanName{DTildeTrainingProtocol} together with the calibration structures of
Parts~\ref{part:dps} and~\ref{part:training}.

\begin{enumerate}
  \item \textbf{$\rho_{\min}$ calibration.} Before training begins, run the
  benchmark correlation study: sample checkpoint-free model outputs, compute
  $(\RewardDT, R_D)$ pairs on a held-aside validation set, and set
  $\rho_{\min}$ from the empirical correlation distribution (e.g.\ its lower
  confidence bound). This discharges the precondition of
  \LeanName{ProxyFaithfulness_decay_detectable}.
  \item \textbf{$\tau$ calibration.} Generate incoherent high-entropy
  baselines (feasibility-passing shuffled/noise outputs), estimate
  $H_q(y_{\mathrm{random}})$, and set $\tau$ strictly above it. This
  discharges \LeanName{B_utility_gate_tightness} and blocks exploitation path
  (iii) of the Structural Goodhart Theorem.
  \item \textbf{Corpus stationarity schedule.} Fix the external corpus update
  schedule before training; policy outputs are never inserted into
  $\mathcal{H}$. This discharges \LeanName{C_corpus_stationarity_condition}.
  \item \textbf{Probe resampling.} Resample $Q_{\mathrm{probe}}$ fresh every
  batch; no probe string persists across steps
  (\LeanName{D_tilde_no_fixed_target}).
  \item \textbf{Rollout and advantage computation.} Two rollouts per prompt;
  advantages quantize to $\pm 1/\sqrt{2}$
  (\LeanName{twoGRPO_advantage_quantization_unbiased}), so the update is a
  signed-rank (DPO-equivalent) step at roughly $1/8$ of the 16-rollout cost.
  \item \textbf{Checkpoint monitoring.} At every checkpoint compute
  $\rho(\RewardDT, R_D)$ on the held-aside validation set; if
  $\rho < \rho_{\min}$, halt and flag a \LeanName{ProxyFaithfulness}
  violation. Report the final correlation as a first-class result.
\end{enumerate}

\section{Known Gaps Registry}
\label{app:known-gaps}

Central, auditable list of every statement the Lean development assumes rather
than proves, with its promotion path. All other results cited in this document
are machine-checked theorems with no \texttt{sorry}.

\begin{enumerate}
  \item \textbf{KG-1 (analytic imports).} \LeanName{klDivergence_nonneg}
  (Gibbs inequality) and \LeanName{kl_cross_entropy_decomposition}
  (cross-entropy link) in \texttt{Creativity/Probability/KLDivergence.lean}.
  Promotion path: derive from Jensen on the convex function $x\mapsto x\log x$
  over the support, mirroring the promotion of
  \LeanName{entropy_mixture_concavity} (PROOF-09). Note that
  \LeanName{maxEntropy_finite_alphabet} is already \emph{derived from} KG-1
  rather than assumed separately.
  \item \textbf{KG-2 (statistical import).} \LeanName{plugIn_bias_axiom}
  (Miller--Madow bias bound, PROOF-12). Promotion path: Hoeffding's
  inequality plus Lipschitz continuity of $-x\log x$; blocked on a Lean~4
  concentration-of-measure dependency.
  \item \textbf{KG-3 (protocol certificates).}
  \LeanName{D_tilde_no_fixed_target} (probe resampling) and
  \LeanName{B_utility_gate_tightness} ($\tau$ calibration) are certificates
  carried by protocol structures: they are discharged by \emph{running the
  protocol correctly} (Appendix~\ref{app:2grpo-protocol}), not by
  mathematics. Their violation consequences are machine-checked by the
  Structural Goodhart Theorem (PROOF-13).
  \item \textbf{KG-4 (resolved).} Formerly registered gaps now closed:
  \LeanName{entropy_mixture_concavity_axiom} (promoted, PROOF-09),
  \LeanName{productKernel_joint_entropy_eq_sum} (promoted, PROOF-10), and
  \LeanName{gain_zero_iff_productKernel} (removed as incorrect and replaced
  by \LeanName{gain_zero_iff_joint_eq_max} plus
  \LeanName{productKernel_additivity_gain_pos}, PROOF-11).
  \item \textbf{KG-5 (optional shifted threshold).} $f$ for the calibration
  variant \LeanName{CUE_shifted} / \LeanName{Rcreativity_shifted} --- not
  part of the canonical definition~\eqref{eq:Rcreativity}. Empirically
  calibrated per domain; all shifted theorems hold for every real $f$.
  \item \textbf{KG-6 (Benchmark Score hyperparameters).}
  $\alpha$ (expansion-bonus weight; also absorbs bits/nats residual scale),
  $\delta_D$ (D-gate threshold), and $\lambda_G$ (interaction-gain weight)
  are empirically calibrated with no canonical closed form. Status:
  registered assumptions, not proved.
\end{enumerate}

\end{document}
